\documentclass[11pt]{article}

% ===== 基本宏包：尽量保持简单 =====
\usepackage[margin=1in]{geometry}
\usepackage{graphicx}
\usepackage{float}
\usepackage{chngcntr}
\usepackage{booktabs}
\usepackage{multirow}
\usepackage{longtable}
\usepackage{array}
\usepackage{calc}
\usepackage{amsmath, amssymb}
\setlength{\LTcapwidth}{\textwidth}
\usepackage{xurl}
\usepackage{hyperref}
\usepackage[numbers,sort&compress]{natbib}
\usepackage{setspace}
\usepackage{fontspec}
\usepackage{xeCJK}
\providecommand{\tightlist}{\setlength{\itemsep}{0pt}\setlength{\parskip}{0pt}}

% ===== 图表编号风格：全文连续编号，使用 Figure / Table =====
\renewcommand{\figurename}{Figure}
\renewcommand{\tablename}{Table}

% ===== 字体与换行策略 =====
\defaultfontfeatures{Ligatures=TeX,Scale=MatchLowercase}
\IfFontExistsTF{Noto Serif}{\setmainfont{Noto Serif}}{\IfFontExistsTF{DejaVu Serif}{\setmainfont{DejaVu Serif}}{\setmainfont{Latin Modern Roman}}}
\IfFontExistsTF{Noto Sans}{\setsansfont{Noto Sans}}{\IfFontExistsTF{DejaVu Sans}{\setsansfont{DejaVu Sans}}{\setsansfont{Latin Modern Sans}}}
\IfFontExistsTF{Noto Sans Mono}{\setmonofont{Noto Sans Mono}}{\IfFontExistsTF{DejaVu Sans Mono}{\setmonofont{DejaVu Sans Mono}}{\setmonofont{Latin Modern Mono}}}
\IfFontExistsTF{Noto Serif CJK SC}{\setCJKmainfont[AutoFakeBold=3,AutoFakeSlant=.2]{Noto Serif CJK SC}}{\IfFontExistsTF{Source Han Serif SC}{\setCJKmainfont[AutoFakeBold=3,AutoFakeSlant=.2]{Source Han Serif SC}}{\IfFontExistsTF{AR PL UMing CN}{\setCJKmainfont[AutoFakeBold=3,AutoFakeSlant=.2]{AR PL UMing CN}}{\setCJKmainfont[AutoFakeBold=3,AutoFakeSlant=.2]{WenQuanYi Zen Hei}}}}
\IfFontExistsTF{Noto Sans CJK SC}{\setCJKsansfont[AutoFakeBold=3,AutoFakeSlant=.2]{Noto Sans CJK SC}}{\IfFontExistsTF{Source Han Sans SC}{\setCJKsansfont[AutoFakeBold=3,AutoFakeSlant=.2]{Source Han Sans SC}}{\IfFontExistsTF{Droid Sans Fallback}{\setCJKsansfont[AutoFakeBold=3,AutoFakeSlant=.2]{Droid Sans Fallback}}{\setCJKsansfont[AutoFakeBold=3,AutoFakeSlant=.2]{AR PL UMing CN}}}}
\IfFontExistsTF{Noto Sans Mono CJK SC}{\setCJKmonofont[AutoFakeBold=3]{Noto Sans Mono CJK SC}}{\IfFontExistsTF{Droid Sans Fallback}{\setCJKmonofont[AutoFakeBold=3]{Droid Sans Fallback}}{\setCJKmonofont[AutoFakeBold=3]{AR PL UMing CN}}}
\setlength{\emergencystretch}{6em}
\sloppy

% 可选：让行距略宽一些，便于审稿阅读
\onehalfspacing

\title{\textbf{CuraView: 基于GraphRAG增强知识验证的医疗幻觉检测多智能体框架}}

\author{{\small
\begin{tabular}{@{}cccc@{}}
\parbox[t]{0.2200\textwidth}{\centering \textbf{Severin Ye} \\ \textit{School of Computer Science and Engineering} \\ \textit{Kyungpook National University} \\ Daegu, Republic of Korea \\ {\ttfamily\small 6severin9\allowbreak@gmail.\allowbreak{}com}} & \parbox[t]{0.2200\textwidth}{\centering \textbf{Xiao Kong} \\ \textit{School of Computer Science and Engineering} \\ \textit{Kyungpook National University} \\ Daegu, Republic of Korea \\ {\ttfamily\small kongx7546\allowbreak@knu.\allowbreak{}ac.\allowbreak{}kr}} & \parbox[t]{0.2200\textwidth}{\centering \textbf{Xiaopeng He} \\ \textit{West China School of Medicine} \\ \textit{Sichuan University} \\ Chengdu, China \\ {\ttfamily\small hxp19180945016\allowbreak@stu.\allowbreak{}wchscu.\allowbreak{}cn}} & \parbox[t]{0.2200\textwidth}{\centering \textbf{Guangsu Yan} \\ \textit{School of Computer Science and Technology} \\ \textit{Shandong University} \\ Jinan, China \\ {\ttfamily\small 202300800569\allowbreak@mail.\allowbreak{}sdu.\allowbreak{}edu.\allowbreak{}cn}} \\
\end{tabular}
}}

\date{} % Nature 系列投稿一般不需要显示日期

\begin{document}
\maketitle

% =========================
% Abstract
% =========================
\begin{abstract}
临床出院小结等摘要需要从冗长的电子病历（EHR）中提炼关键信息，人工成本高。大语言模型可提升生成效率，但易产生与病历不一致的``忠实性幻觉''，在医疗场景中带来安全风险。为此，本文提出 \textbf{CuraView}：以患者级 EHR 构建 GraphRAG 知识图谱，并在生成---检测闭环中对句子级陈述进行证据检索与分级（E1--E4），输出结构化、可解释的支持/冲突证据链。我们在 Discharge-Me benchmark 的 250 名患者子集上评估（50 名测试）。基于 Qwen3-14B 的微调检测模型在安全关键的 E4 指标上达到 \textbf{F1=0.831}（召回率 \textbf{90.9\%}、精确率 \textbf{76.5\%}），在 E3+E4 指标上达到 \textbf{F1=0.823}；其在 E4 检测上相较基座模型实现了 \textbf{50.0\%} 的相对提升，并优于 RAGTruth 风格和 QAGS 风格基线。结果表明，基于患者证据链的图检索核验能够有效提升临床摘要的事实可靠性，并生成可复用的标注数据，用于后续训练与蒸馏。

关键词---临床摘要；事实一致性；幻觉检测；GraphRAG；知识图谱
\end{abstract}

% =========================
% Introduction
% =========================
\section{引言}
大语言模型(Large Language Models, LLMs)在各个领域的人工智能应用中引发了革命,在医疗保健领域展现出尤为变革性的潜力。近期的模型如GPT-4、Med-PaLM 2及其临床变体在医学问答方面表现出显著能力。其中,Med-PaLM 是首个在 MedQA 的 USMLE 风格题目上超过及格线（\textgreater60\%）的AI系统 \citep{singhal2023clinicalKnowledge},Med-PaLM 2 则进一步将该基准准确率提升至 86.5\% \citep{singhal2023expertMedicalQA}; 与此同时,GPT-4 在 USMLE 官方练习材料和 MultiMedQA 上也展现出强竞争力 \citep{nori2023gpt4Medical}。更重要的是，LLM 的应用已经从 benchmark 问答扩展到临床工作流相关任务，包括诊断辅助、临床文档编写、放射学报告生成、出院小结生成和患者教育 \citep{lee2023gpt4Medicine,thirunavukarasu2023llmMedicine,xu2024dischargeMeSharedTask}。其中，出院小结生成已进一步发展为统一 shared task \citep{xu2024dischargeMeSharedTask}，而关于医生时间分配与文档负担的研究则表明，临床文档自动化确实对应明确的现实需求 \citep{sinsky2016physicianTime,zhao2025clinicalDocumentationBurden}。市场预测也显示医疗 AI 具有显著经济空间 \citep{marketsandmarkets2025healthcareAI}，但经济潜力并不会自动转化为临床落地能力。当前医疗AI面临的核心障碍,是由幻觉引发的``临床信任鸿沟'':如果模型生成内容无法被患者级证据稳定核验,医生和机构就难以在真实工作流中依赖它。

然而,这一光明前景面临着一个关键障碍:\emph{幻觉(hallucinations)}------生成看似合理但实际上错误的信息。虽然通用领域LLM的幻觉可能只造成不便,但医疗幻觉却对患者安全和临床信任构成严重风险 \citep{alkaissi2023chatgptHallucinations}。已有综述和实证研究表明，医疗场景中的 LLM 可能在问答、摘要和临床支持任务中生成看似流畅但事实上错误的建议 \citep{thirunavukarasu2023llmMedicine,rajpurkar2022aiHealthMedicine}。医疗LLM中的幻觉表现为自信陈述但实际上不正确的医疗信息，范围从误诊到虚构的检查结果。基于对临床文档错误模式的系统分析，我们识别出七种关键的幻觉类型，覆盖诊断、用药、检验结果、时间、数值、否定和虚构事实等错误模式（详细定义见第3.2节）。医疗AI的风险尤为严重。单个用药剂量错误可能致命;遗漏癌症诊断可能延误挽救生命的治疗;虚构的检查结果可能引发不必要的干预。现有研究还表明,医疗文本生成中的幻觉具有明显的任务依赖性,不同任务中的错误形态、评估粒度和风险负担并不一致,因此并不存在一个统一的``医疗LLM幻觉率''指标。本文将更细致的任务差异与代表性文献留待第2.1节统一说明,这里仅强调一点：一旦错误进入临床文档生成场景,其风险将直接转化为患者安全问题,因而需要面向患者级事实证据的专门检测机制。

目前的医疗幻觉研究存在四个关键局限性，这也是我们工作的动力。\textbf{L1：缺乏系统性生成。}现有的医疗幻觉基准仍高度依赖人工策展和专家标注，例如 Med-HALT \citep{pal2023medHalt}，成本高昂且规模有限。此前尚无工作提供一个能够大规模系统性创建多样化医疗错误的\textbf{自动化}幻觉生成框架。\textbf{L2：缺失患者特定上下文。}大多数研究仍在通用医学问题、跨患者汇总语料或通用临床生成任务上评估幻觉,而没有把来自电子健康档案（EHR）的\textbf{个体化患者数据}直接组织为可核验的事实依据 \citep{rotmensch2017healthKnowledgeGraph,finlayson2014graphMedicine}。例如，检测''患者服用阿司匹林''是否为幻觉，需要了解患者的过敏史和当前用药情况------而现有基准测试通常缺乏这种患者级上下文。\textbf{L3：证据支撑不足。}现有医疗幻觉研究大多报告题目级正确率、句子级幻觉率或文书级错误负担，但较少提供面向患者级核验的结构化解释，说明\textbf{为什么}某个陈述是错误的，或者\textbf{什么证据}与其矛盾 \citep{tonekaboni2019cliniciansWant,asgari2025clinicalSafetySummarisation,williams2025dischargeSummaries}。这限制了可解释检测方法的发展。\textbf{L4：仅关注评估。}之前的工作重点在于\textbf{评估}现有模型的幻觉率 \citep{pal2023medHalt}，但没有提供在生产环境中生成、检测和解释幻觉的\textbf{端到端系统}。这四项局限性并非彼此孤立，而是分别导向本文的四个研究问题：L1 对应 RQ1 的系统化幻觉生成，L2 对应 RQ2 的患者特定知识图谱构建，L3 对应 RQ3 的证据增强检测，L4 对应 RQ4 的端到端框架设计。

为了构建临床可行的医疗幻觉检测系统，我们必须解决四个基本研究问题：\textbf{RQ1：多样化医疗幻觉的系统化生成。}现有的医疗幻觉检测基准在规模和多样性方面都有限 \citep{pal2023medHalt}。人工标注幻觉成本高昂，需要专家临床医生识别数千名患者中的细微错误。此外，真实世界的医疗错误表现出多样化的模式------从明显的事实矛盾到细微的临床不合理性------需要系统化覆盖。\textbf{RQ2：患者特定知识图谱构建。}电子健康档案包含跨异构表(诊断、用药、检查结果、生命体征、临床笔记)的丰富患者信息，但将其转换为适合幻觉验证的结构化知识并非易事。\textbf{RQ3：基于知识图谱的上下文增强检测。}给定患者的知识图谱，检测幻觉需要对复杂的实体关系和时间约束进行推理。\textbf{RQ4：端到端框架设计。}一个完整的幻觉检测系统需要在具有不同技术需求的多个组件之间无缝协调。这些研究问题带来了相应的技术挑战：\textbf{C1：}针对RQ1，我们需要一个自动化幻觉生成框架，能够：(1)创建反映现实临床错误的多样化错误类型；(2)保持医疗合理性以避免平凡检测；(3)保留句子级粒度以进行细粒度评估；(4)提供带有基于证据解释的真实标签。\textbf{C2：}针对RQ2，我们需要一个知识图谱构建流程，能够：(1)从多表EHR数据中提取医疗实体和关系；(2)处理临床术语变化和缩写；(3)捕捉患者状况的时间演变；(4)在保护患者隐私的同时实现跨患者知识聚合；(5)扩展到数万名患者。\textbf{C3：}针对RQ3，我们需要一个检测机制，能够：(1)高效检索相关知识图谱上下文(避免对大型图谱进行穷举搜索)；(2)执行多跳推理以验证复杂的临床陈述；(3)以置信度分数区分支持、不支持和矛盾的陈述；(4)生成引用具体EHR证据的人类可解释解释。\textbf{C4：}针对RQ4，我们需要一个系统架构，能够：(1)编排幻觉生成和检测代理，使用一致的数据格式；(2)支持基于API和本地模型部署以实现灵活性；(3)实现对大型患者群体的高效批处理；(4)实施稳健验证以确保生成-检测一致性；(5)提供全面的评估指标，比较不同配置下的系统性能。

为了应对这些挑战，我们提出\textbf{CuraView}，一个具有GraphRAG增强知识验证的医疗幻觉检测多智能体框架。这里采用 GraphRAG 而非 vanilla RAG 的原因在于，临床核验并不只是``找到相似文本''，而是要在诊断、用药、检验结果和时间信息之间恢复患者特定的关系结构。GraphRAG 能够把多表 EHR 中分散的证据组织为可查询的关系图谱，因此更适合支持句子级事实核验。CuraView 通过以下方式解决了上述四个局限性：针对\textbf{L1}，提供具有七种错误类型的自动化幻觉生成框架；针对\textbf{L2}，构建基于EHR的患者特定知识图谱；针对\textbf{L3}，实施证据分级标注（E1-E4）以提供可解释的检测结果；针对\textbf{L4}，提供完整的生成-检测-评估端到端流程。我们的工作做出以下关键贡献：

\begin{itemize}
\tightlist
\item
  贡献1：临床文档导向的医疗幻觉分类法
\item
  贡献2：端到端多智能体框架
\item
  贡献3：GraphRAG驱动的医疗知识图谱
\item
  贡献4：统一的结构化输出可靠性流水线
\item
  贡献5：全面的评估框架
\end{itemize}

% =========================
% 相关工作
% =========================
\section{相关工作}
在详细介绍CuraView方法之前，我们首先综述医疗幻觉检测、知识图谱增强生成和多智能体系统的现有研究，以建立本研究的定位和创新点。

\subsection{医疗大语言模型中的幻觉}

\subsubsection{医疗大语言模型的能力与局限性}

近年来，将大语言模型（LLM）应用于医疗领域取得了显著进展。基础能力方面，\textbf{Med-PaLM} \citep{singhal2023clinicalKnowledge} 首次在 MedQA 的 USMLE 风格题目上超过及格线，\textbf{Med-PaLM 2} \citep{singhal2023expertMedicalQA} 则进一步把多项医学问答基准推进到当时的最先进水平。与此同时，\textbf{GPT-4} 在 USMLE 官方练习材料和 MultiMedQA 上也表现出强竞争力 \citep{nori2023gpt4Medical}。在领域适配方向，\textbf{Clinical Camel} \citep{toma2023clinicalCamel} 和 \textbf{PMC-LLaMA} \citep{wu2023pmcLLaMA} 展示了基于医学语料继续训练的模型在临床笔记生成与医学问答中的有效性。换言之，现有工作已经证明医疗 LLM 可以``会答题、会生成''，但这并不等同于``可在高风险文档场景中稳定可信地生成''。

然而，这些令人印象深刻的基准测试表现掩盖了一个关键弱点：\emph{幻觉}。Ji 等人 \citep{ji2023hallucinationSurvey} 对自然语言处理中的幻觉进行了全面综述，将其分为\textbf{事实性幻觉}（与既定知识相矛盾）和\textbf{忠实性幻觉}（与源文档不一致）。在医疗背景下，这两类幻觉都会带来严重的风险。

\subsubsection{医疗幻觉基准测试}

\textbf{Med-HALT} \citep{pal2023medHalt} 引入了第一个系统性的医疗幻觉基准，涵盖了五个医学专业的 1,755 个问题。他们发现，即使是像 GPT-4 这样最先进的模型，在事实性医学问题上的幻觉率也达到了 10-15\%。至关重要的是，幻觉往往伴随着高置信度评分，这使得它们在临床环境中特别危险。

进一步地，现实世界中也已经出现过 AI 生成医疗建议包含危险幻觉的公开案例 \citep{pal2023medHalt,ji2023hallucinationSurvey}，这说明相关风险并不只存在于 benchmark 环境，而是会直接转化为对可靠检测机制的现实需求。

\textbf{Asgari等} \citep{asgari2025clinicalSafetySummarisation} 将评估粒度从``题目级正确率''推进到\textbf{句子级临床摘要安全性评估}。在12,999个经临床医生标注的句子中，他们报告了1.47\%的句子级幻觉率和3.45\%的遗漏率，其中44\%的幻觉被归类为重大错误。这说明在相对结构化的临床笔记摘要任务中，绝对幻觉率可能较低，但一旦出现错误，其临床风险并不低。

\textbf{Williams等} \citep{williams2025dischargeSummaries} 则考察了\textbf{出院小结叙述生成}。该研究没有将问题压缩为单一幻觉率，而是从不准确、遗漏和幻觉三类错误出发，比较每份文书中的独特错误总数。结果显示，LLM生成版本平均包含2.91个独特错误，高于医生撰写版本的1.82个。这提示出院小结场景中更适合采用``每份文书错误负担''而非单一句子级幻觉率来衡量风险。

近期的 \textbf{Discharge Me} 共享任务进一步把出院小结生成从单篇案例研究推进为统一 benchmark \citep{xu2024dischargeMeSharedTask}。在该 benchmark 上，\textbf{EPFL-MAKE} 通过上下文窗口扩展与动态信息选择处理超长病历 \citep{wu2024epflMakeDischargeMe}，而 \textbf{UF-HOBI} 则采用``抽取关键临床概念 + 生成连贯段落''的两阶段混合流水线 \citep{lyu2024ufHobiDischargeMe}。这些工作说明，出院小结生成已经形成相对清晰的任务谱系，但系统目标仍主要是``生成更像样的文书''，而不是对患者级事实进行逐句核验并给出证据化解释。

在\textbf{放射学报告生成}中，风险表现得更为突出。\textbf{Artsi等} \citep{artsi2025radiologyReview} 的系统综述纳入了15项研究，发现所有模型都出现了不同程度的hallucinations、misdiagnoses或报告不一致问题；进一步地，综述指出缺失临床细节在12/15项研究中被报告，显示该领域的错误不仅涉及``虚构''，也涉及``遗漏''和``错配''。

\textbf{RadFlag} \citep{zhang2024radFlag} 从句子级检测角度对放射报告幻觉进行了更细粒度分析。该研究指出，高性能放射报告生成模型中约40\%的生成句子可能包含幻觉；在其报告级评估中，被系统标记为高风险的报告平均包含4.2个真实幻觉，而被接受的报告平均为1.9个。这进一步说明放射报告任务中的幻觉密度和报告级风险都可能显著高于其他更结构化的医疗文本生成场景。

综上，现有文献更支持一种\textbf{任务依赖、评估粒度依赖}的理解：医学问答、临床摘要、出院小结和放射报告生成中的``幻觉''并不是同一个可直接横向比较的指标。本研究在第3.2节提出的细粒度分类法，专注于临床文档中的句子级错误，并强调与患者特定EHR证据的对齐，从而与此前聚焦问答或放射摘要的工作形成互补。

\subsection{知识图谱与检索增强生成}

\subsubsection{检索增强生成基础}

检索增强生成（RAG） \citep{lewis2020retrievalAugmentedGeneration} 通过将神经检索与语言生成相结合，引入了 LLM 应用的范式转变。经典的 RAG 架构使用密集向量检索（DPR）从语料库中获取相关文档，然后使用像 BART 这样的模型根据检索到的上下文进行生成。这种方法显著提高了在知识密集型任务（如开放域问答）上的性能。

随后的工作将 RAG 扩展到事实核验与领域问答。\textbf{RARR} \citep{gao2023rarr} 的重点是用外部证据对模型输出进行检索后修订；在医学场景中，\textbf{MedRAG} \citep{xiong2024medRAG} 系统比较了不同医学 RAG 配置；在生物领域，\textbf{BioRAG} \citep{wang2024bioRAG} 则结合科学文献库、领域特定嵌入和知识层级来支持复杂问题推理。尽管这些工作表明检索增强在专业领域具有明显价值，但其核心仍主要建立在\textbf{扁平化文档检索}之上，尚不足以显式建模医疗概念之间的结构化关系。例如，在出院小结核验中，系统不仅要找到``患者正在使用某药物''的文本，还要将其与其他记录中的过敏史、停药记录或异常检验结果关联起来。扁平检索通常只能返回局部相似片段，却难以稳定恢复这种跨来源关系。

\subsubsection{GraphRAG：图结构检索增强}

微软的 \textbf{GraphRAG} \citep{edge2024graphRAGSummarization} 通过在检索前将文档转换为知识图谱，解决了传统 RAG 对文档结构化关系建模不足的问题。该方法采用两级索引架构：实体级索引（局部搜索）用于检索特定实体及其直接关系，而社区级索引（全局搜索）则从 Leiden 算法检测出的社区中聚合知识，提供更广泛的上下文理解。

GraphRAG 的图构建流程包括四个关键步骤：首先利用大语言模型从文本中提取实体和关系，构建初始知识图谱；其次应用 Leiden 社区检测算法识别语义相关的实体簇；随后为每个社区生成摘要性描述；最后构建结合结构化图谱知识和非结构化文本的混合向量索引。相较于标准 RAG，这种图结构方法能够捕捉复杂的多跳推理关系，并通过社区结构发现文档集合中的潜在主题模式。

总体而言，现有研究呈现出从无显式外部证据、到扁平检索增强、再到图结构检索增强的技术演进脉络。在医疗场景中，近期 \textbf{Medical Graph RAG} 工作进一步表明，将患者文档连接到可信医学来源、并结合图检索与响应细化，可以显著提升证据性与可信度 \citep{wu2025medicalGraphRAG}；相应地，医疗 RAG 系统综述也指出，外部知识源是缓解医疗幻觉与提升透明度的重要技术路径 \citep{amugongo2025healthcareRAGReview}。在需要患者特定证据核验、多跳关系推理和跨来源信息整合的医疗幻觉检测场景中，这一脉络也解释了为什么 GraphRAG 比无检索或仅基于扁平文本片段的检索方式更适合作为本文的方法基础。

\subsubsection{医学知识图谱}

医学本体和知识图谱提供了临床知识的结构化表示：

\begin{itemize}
\item
  \textbf{UMLS（统一医学语言系统）} \citep{bodenreider2004umls}：整合了 200 多种医学术语集，实现了语义互操作性，但缺乏患者特定的信息
\item
  \textbf{SNOMED-CT} \citep{donnelly2006snomedCt}：具有层级关系的全面临床术语系统
\item
  \textbf{ICD 编码} \citep{who2016icd10}：国际疾病分类标准
\end{itemize}

\textbf{基于 EHR 的知识图谱构建：} Rotmensch 等人 \citep{rotmensch2017healthKnowledgeGraph} 从 EHR 数据中自动构建疾病知识图谱，通过 ICD 编码和临床笔记发现新的疾病关联。Finlayson 等人 \citep{finlayson2014graphMedicine} 为临床笔记开发了实体链接方法，将提及内容与 UMLS 概念对齐。

\textbf{医学知识图谱构建中的挑战：}

\begin{enumerate}
\def\labelenumi{\arabic{enumi}.}
\item
  \textbf{术语变体}：临床笔记中的缩写、俗语、拼写错误
\item
  \textbf{时序建模}：疾病进展、治疗历史的演变
\item
  \textbf{隐私保护}：去标识化与知识完整性之间的权衡
\item
  \textbf{跨患者聚合}：平衡知识共享与隐私保护
\end{enumerate}

除通用本体资源外，近期关于 \textbf{SNOMED CT 与大语言模型结合} 的综述还显示，术语知识可以通过输入增强、额外融合模块或检索式知识注入进入下游 NLP 与 NLG 流程，并在不少任务上带来性能改进 \citep{chang2024snomedCtLLM}。这也为本文后续将结构化医疗术语与患者级图谱结合提供了方法论依据。

\subsubsection{我们的方法定位}

\textbf{关键区别：}

\begin{enumerate}
\def\labelenumi{\arabic{enumi}.}
\item
  \textbf{患者特定图谱}：面向单个患者构建个体化图谱，而非依赖通用医学知识库
\item
  \textbf{面向核验的证据语义}：采用 E1-E4 证据分级支撑句子级事实核验
\item
  \textbf{多表 EHR 集成}：统一处理诊断、用药、检验结果和生命体征等异构来源
\item
  \textbf{动态患者证据组织}：围绕当前患者的真实记录进行结构化建模，而非使用静态通用本体
\end{enumerate}

\subsection{多智能体系统与 LLM 应用}

\subsubsection{LLM 智能体框架}

\textbf{LangChain} \citep{chase2023langchain} 提供了一个模块化框架，包括提示词模板、LLM 封装器、工具调用接口和记忆管理。我们基于 LangChain 1.0 构建了 CuraView 的生成和检测智能体，以确保可扩展性。

\textbf{AutoGPT} \citep{significantgravitas2023autogpt} 展示了具有长期记忆和互联网访问权限的自主目标分解能力。\textbf{ReAct} \citep{yao2023react} 引入了''推理+行动''范式，在迭代推理中交替使用思维链与工具调用。

近期 agent 文献还表明，参数规模并不是工具型智能体性能的唯一决定因素。在 function calling、多轮 tool use 与 workflow orchestration 等结构化子任务上，经过高质量 agent 数据合成、监督微调或规则奖励强化学习优化的小模型，可以达到与更大通用模型相当、甚至更优的表现 \citep{liu2025toolACE,prabhakar2025apiGenMt,zhang2025nemotronToolN1}。其中，\textbf{ToolACE} \citep{liu2025toolACE} 关注 function-calling 能力，\textbf{APIGen-MT} \citep{prabhakar2025apiGenMt} 强调多轮 agent 交互，而 \textbf{Nemotron-Research-Tool-N1} \citep{zhang2025nemotronToolN1} 则展示了规则奖励强化学习对工具调用模型的提升作用。它们共同支持一个判断：在 agent 场景中，任务特化、结构化接口和训练数据质量在某些子任务上可以比单纯扩大参数规模更关键。

\subsubsection{与现有工作的比较}

现有的医疗多智能体系统中，\textbf{MedAgents} \citep{tang2023medAgents} 这类代表性工作主要聚焦于诊断与治疗任务，依赖通用医学知识库（PubMed、教材等）进行推理，但缺乏针对LLM生成内容质量的验证机制。这些系统在处理患者特定场景时面临两个关键局限：一是无法利用个体化的电子健康档案（EHR）进行精准推理，二是对生成内容的幻觉问题缺少系统性检测手段。从方法论上看，CuraView 的双智能体设计结合了两条正在形成的技术脉络。一方面，对抗式样本生成能够持续构造更具迷惑性的错误，从而提升检测模块的鲁棒性；另一方面，LLM-as-a-judge 范式表明，围绕证据执行结构化评审比单次开放式回答更适合高风险质量控制场景。基于这两点，生成智能体与检测智能体的解耦并非单纯工程拆分，而是服务于``持续制造挑战样本''和``证据驱动核验判断''这两类互补任务。

CuraView 在以下四个维度实现了关键突破：\textbf{（1）任务定位}------首个专注于幻觉检测与质量控制的医学多智能体系统，而非传统的诊疗辅助；\textbf{（2）知识来源}------构建患者特定的 EHR 知识图谱，提供个体化的事实依据，而非依赖泛化的医学文献；\textbf{（3）智能体协作}------采用对抗式生成-检测架构，以系统化生成样本持续挑战检测智能体的鲁棒性；\textbf{（4）数据贡献}------在该对抗流程中同步沉淀可用于后续训练与蒸馏的数据资源。这里的重点在于与现有工作的差异，具体系统实现见第4节。

进一步地，相较于现有主要停留在提示工程、通用工具调用或单次模型输出层面的医学智能体系统，CuraView 还强调结构化输出的系统可用性，即生成结果能否稳定进入后续验证、评估和批处理流程。这一工程差异是本文与通用医学 agent 框架的重要分野，完整机制说明见第4节。

% =========================
% 问题定义：医疗幻觉的定义、分类与证据分级
% =========================
\section{问题定义：医疗幻觉的定义、分类与证据分级}
\subsection{幻觉定义}

在大语言模型（LLMs）的背景下，\emph{幻觉}指模型生成的内容看似合理但实际上与真实知识或源文档相矛盾的现象 \citep{ji2023hallucinationSurvey}。Ji等人将幻觉分为两类：\textbf{事实性幻觉}（与既定知识相矛盾）和\textbf{忠实性幻觉}（与源文档不一致）。

本研究聚焦于医疗领域的\textbf{忠实性幻觉}，特别是在临床文档生成过程中产生的错误医疗信息。与通用领域不同，医疗幻觉对患者安全构成直接威胁------单个用药剂量错误可能致命，遗漏癌症诊断可能延误挽救生命的治疗，虚构的检查结果可能引发不必要的干预。

如引言所示，现有文献已经在临床摘要、出院小结和放射学报告等任务中反复观察到医疗幻觉风险，但不同任务的错误率、严重性和评估粒度差异明显，因此并不存在统一的``医疗LLM幻觉率''指标 \citep{pal2023medHalt,ji2023hallucinationSurvey,asgari2025clinicalSafetySummarisation,williams2025dischargeSummaries,artsi2025radiologyReview,zhang2024radFlag}。本节不再重复展开各任务统计，而转向本文后续方法直接依赖的定义性内容，即临床文档场景下的幻觉分类法与证据分级体系。

\subsection{医疗幻觉分类法}

基于对临床文档错误模式的系统分析，我们将医疗幻觉分为七种类型，每种类型对应不同的临床风险级别：

\begin{enumerate}
\def\labelenumi{\arabic{enumi}.}
\item
  \textbf{诊断错误（Diagnostic Errors）}：疾病识别错误，如将"肺炎"误写为"肺结核"。这类错误可能导致完全错误的治疗方案，具有极高的临床风险。
\item
  \textbf{用药错误（Medication Errors）}：药物或剂量错误，如将"阿司匹林81mg"误写为"阿司匹林325mg"。剂量错误可能导致药物中毒或治疗无效，是患者安全的重要威胁。
\item
  \textbf{检验结果错误（Test Result Errors）}：虚构或错误的检查数值，如将"血糖96mg/dL"误写为"血糖196mg/dL"。这类错误可能触发不必要的紧急干预或延误必要的治疗。
\item
  \textbf{时间错误（Temporal Errors）}：时间信息不一致，如将"昨天入院"误写为"上周入院"。时间错误影响对病程的理解和治疗时机判断。
\item
  \textbf{数值错误（Numerical Errors）}：生命体征数值错误，如将"血压120/80mmHg"误写为"血压180/120mmHg"。这类错误可能引发错误的紧急响应或掩盖真实的危急情况。
\item
  \textbf{否定错误（Negation Errors）}：肯定/否定关系颠倒，如将"无糖尿病史"误写为"有糖尿病史"。看似微小的否定词变化可能完全改变患者的医疗风险评估。
\item
  \textbf{虚构事实（Fabricated Facts）}：完全虚构的临床事件，如添加不存在的手术记录。这类错误通常在证据中没有被提及，但根据模型对上下文和临床常识的判断应视为可疑或错误，是典型的 E3 类幻觉。
\end{enumerate}

该分类法与之前的工作 \citep{rotmensch2017healthKnowledgeGraph,ji2023hallucinationSurvey}不同，专注于\textbf{临床文档}中的\textbf{句子级}错误，从而实现细粒度的评估和定位。

\subsection{证据分级系统}

为了量化幻觉的可检测性和临床影响，我们引入了四级证据分级系统（E1-E4）：

\begin{itemize}
\item
  \textbf{E1（强支持）}：陈述有明确的EHR证据支持，如患者记录中存在相应的诊断代码、用药记录或检查结果。
\item
  \textbf{E2（弱支持）}：陈述有间接证据支持，如可以从相关信息推断，但无直接记录。例如，即使病历中缺少明确的糖尿病诊断代码，若患者长期使用二甲双胍等降糖药物，则``患者存在糖尿病''可视为一种典型的 E2（弱支持）情形。
\item
  \textbf{E3（无支持）}：陈述在EHR证据中没有被提及，也缺乏直接支持，但根据模型对上下文的判断仍应视为有问题；典型情形包括无中生有、凭空添加事件或加入病历中不存在的事实。
\item
  \textbf{E4（矛盾）}：陈述与EHR中的明确事实证据直接矛盾，如诊断、用药、化验值或否定关系与记录不一致。这是最关键的安全等级，代表可被证据明确判定的错误信息。
\end{itemize}

本研究重点关注\textbf{E3（无支持）}和\textbf{E4（矛盾）}级幻觉的检测，因为这两类代表了需要系统性验证的临床陈述。其中，E4级幻觉是本文主要关注的安全关键错误，而E3级幻觉则代表证据中未提及、但根据模型判断应被视为有问题的可疑陈述。

% =========================
% 方法
% =========================
\section{方法}
\subsection{系统架构概览}

CuraView 是一个端到端的医学幻觉检测框架，由三个核心组件协同工作。

幻觉生成智能体基于 LangChain 1.0 构建，系统性地生成多样化的医学错误。

GraphRAG 知识图谱从电子健康记录（EHR）中构建患者特定的知识库，支持上下文增强检索。

幻觉检测智能体利用知识图谱对生成文本进行验证，并基于证据识别幻觉。这三个组件形成了一个完整的生成-验证-评估闭环系统。

\begin{figure}[H]
\centering
\includegraphics[width=0.92\textwidth]{/home/severin/Codelib/CuraView/论文\_md形式/figures/figure\_01.png}
\caption{CuraView系统架构概览。框架由三个核心组件协同工作：幻觉生成智能体基于LangChain系统化生成医疗错误；GraphRAG知识图谱从EHR构建患者特定知识库；幻觉检测智能体利用知识图谱进行证据验证。}
\label{fig:system-architecture}
\end{figure}

\subsubsection{数据流}

系统通过五个阶段处理数据。

\textbf{数据准备阶段}将MIMIC-IV多表数据进行整合转换，生成统一的患者记录。

\textbf{知识图谱构建阶段}对患者数据执行实体抽取、关系识别和社区检测，最终存储到数据库中。

\textbf{幻觉生成阶段}从原始出院文本开始，通过句子切分、适用性评估和幻觉生成，产生重写文本及相应的幻觉记录。

\textbf{幻觉检测阶段}对重写文本进行句子切分，通过GraphRAG查询获取知识上下文，结合LLM推理生成检测结果。

\textbf{对比评估阶段}整合生成结果和检测结果，计算召回率、精确率和F1分数，生成评估报告。

\subsubsection{技术框架}

采用LangChain框架实现智能体功能,使用GraphRAG进行知识图谱构建与检索。模型部署支持API调用和本地推理两种方式,分别使用大规模商用模型和轻量化开源模型。文本处理采用SaT (Segment any Text)模型\citep{frohmann2024segmentAnyText}配合规则方法进行句子切分,图构建使用社区检测算法进行层级组织。

\subsection{幻觉生成智能体}

\subsubsection{设计原则}

幻觉生成智能体旨在在保持医学合理性的同时，系统性地生成多样化的医学错误模式。我们确立了四项核心原则：

\textbf{P1：医学合理性} ------ 仅对包含医疗信息的句子进行幻觉生成，生成的幻觉必须在医学上合理，以避免被轻易检测（例如，避免明显不可能的数值，如"血压 500/300"）。

\textbf{P2：类型多样性} ------ 覆盖七种临床错误类型，每一种都有独立的生成策略。

\textbf{P3：可控采样} ------ 通过参数配置重写比例（默认 40\%），在数据多样性与标注成本之间取得平衡。

\textbf{P4：证据可追溯性} ------ 每一个幻觉都标注证据等级（E1--E4），用于后续检测评估。

除了支持幻觉检测这一核心功能外，该智能体还具有重要的数据合成价值：系统化的幻觉生成流程能够持续构建高质量医学幻觉标注数据，每条样本至少包含原始文本、幻觉文本、幻觉类型、证据等级和详细解释；与此同时，检测阶段产生的推理轨迹和知识检索记录也可进一步沉淀为蒸馏数据，用于将大模型的推理能力迁移到轻量化模型中。也就是说，生成与检测不仅构成质量控制闭环，同时也构成后续模型迭代的数据来源。

\subsubsection{两阶段句子过滤机制}

为确保生成质量，我们设计了一个两阶段句子过滤流程：

\textbf{阶段1：句子适用性评估。}系统首先将出院小结文本切分为独立句子，然后由LLM逐一评估每个句子是否适合进行幻觉改写。评估的三个核心判据为：（1）句子是否包含可验证的医学事实（如诊断、用药、检验数值），而非患者基本信息或主观描述；（2）改写后是否能保持医学合理性，避免生成明显不可能的内容；（3）句子复杂度是否适中，既不过于简单（如"患者入院"）也不过于复杂导致难以控制改写质量。系统采用批量并发处理提高效率，并在API调用失败时自动重试。

\textbf{阶段2：幻觉生成。}对于通过句子适用性评估的句子，系统执行以下步骤：首先，从患者的完整EHR数据中提取与该句子相关的证据信息，包括诊断记录、用药信息、检验结果、生命体征等；其次，LLM基于提取的证据和原始句子，按照七种幻觉类型之一生成改写，确保生成的错误在医学上合理但与证据不符；最后，系统构建结构化输出，包含修改后的幻觉文本、所属幻觉类型（如诊断错误、用药错误等）、修改原因的详细解释、以及基于EHR证据的证据等级标注（E1-E4）。生成结果经过格式和逻辑一致性双重验证，不符合要求的结果将被丢弃或重新生成。

\subsubsection{七种幻觉类型}

Table 1 给出了我们的幻觉分类体系及示例。

\begin{longtable}[]{@{}
  >{\raggedright\arraybackslash}p{(\columnwidth - 8\tabcolsep) * \real{0.2000}}
  >{\raggedright\arraybackslash}p{(\columnwidth - 8\tabcolsep) * \real{0.2000}}
  >{\raggedright\arraybackslash}p{(\columnwidth - 8\tabcolsep) * \real{0.2000}}
  >{\raggedright\arraybackslash}p{(\columnwidth - 8\tabcolsep) * \real{0.2000}}
  >{\raggedright\arraybackslash}p{(\columnwidth - 8\tabcolsep) * \real{0.2000}}@{}}
\caption{七种幻觉类型及示例}\tabularnewline
\toprule\noalign{}
\begin{minipage}[b]{\linewidth}\raggedright
类型
\end{minipage} & \begin{minipage}[b]{\linewidth}\raggedright
定义
\end{minipage} & \begin{minipage}[b]{\linewidth}\raggedright
原始
\end{minipage} & \begin{minipage}[b]{\linewidth}\raggedright
幻觉
\end{minipage} & \begin{minipage}[b]{\linewidth}\raggedright
证据
\end{minipage} \\
\midrule\noalign{}
\endfirsthead
\toprule\noalign{}
\begin{minipage}[b]{\linewidth}\raggedright
类型
\end{minipage} & \begin{minipage}[b]{\linewidth}\raggedright
定义
\end{minipage} & \begin{minipage}[b]{\linewidth}\raggedright
原始
\end{minipage} & \begin{minipage}[b]{\linewidth}\raggedright
幻觉
\end{minipage} & \begin{minipage}[b]{\linewidth}\raggedright
证据
\end{minipage} \\
\midrule\noalign{}
\endhead
\bottomrule\noalign{}
\endlastfoot
diagnosis\_error & 错误诊断 & 诊断为肺炎 & 诊断为肺结核 & E4 \\
medication\_error & 错误药物/剂量 & 阿司匹林 81mg & 阿司匹林 325mg & E4 \\
exam\_result\_error & 虚构化验值 & 血糖 96 mg/dL & 血糖 196 mg/dL & E4 \\
time\_error & 时间错误 & 入院 1/5 & 入院 1/15 & E4 \\
value\_error & 错误生命体征 & 血压 120/80 & 血压 180/120 & E4 \\
negation\_error & 肯定/否定颠倒 & 无糖尿病史 & 有糖尿病史 & E4 \\
invented\_fact & 虚构事件 & - & 接受阑尾切除术 & E3 \\
\end{longtable}

\subsubsection{证据分级系统}

本文采用第3.3节定义的 E1-E4 证据分级体系，并在生成阶段将其作为样本标注协议。在具体实现中，六种与既有 EHR 事实直接冲突的类型（diagnosis\_error、medication\_error、exam\_result\_error、time\_error、value\_error、negation\_error）标注为 E4；invented\_fact 因其在证据中未被提及但仍应视为可疑错误，标注为 E3。对应地，在后续评估中，我们将 E4 作为主要安全关键结果，将 E3+E4 作为补充的宽口径指标。

\subsection{GraphRAG 知识图谱构建}

\subsubsection{实体与关系抽取}

Table 2 总结了我们定义的实体类型及统计信息。系统定义了十种关键关系类型连接实体，包括诊断关系、用药关系、症状关系、操作关系、生命体征关系、检验关系、科室关系以及临床指征关系等。

\begin{longtable}[]{@{}llll@{}}
\caption{实体类型及统计}\tabularnewline
\toprule\noalign{}
实体类型 & 描述 & 每患者平均 & 示例 \\
\midrule\noalign{}
\endfirsthead
\toprule\noalign{}
实体类型 & 描述 & 每患者平均 & 示例 \\
\midrule\noalign{}
\endhead
\bottomrule\noalign{}
\endlastfoot
PATIENT & 患者主体 & 1.0 & ``Patient'' \\
DIAGNOSIS & 诊断 & 3.2 & ``Pneumonia'' \\
MEDICATION & 药物 & 5.7 & ``Aspirin'' \\
LAB\_TEST & 检查类型 & 2.8 & ``CBC'' \\
LAB\_RESULT & 化验结果 & 6.5 & ``Glucose 96'' \\
VITAL\_SIGN & 生命体征 & 4.1 & ``BP 120/80'' \\
SYMPTOM & 症状 & 2.3 & ``Chest pain'' \\
PROCEDURE & 操作 & 1.8 & ``ECG'' \\
DEPARTMENT & 科室 & 1.0 & ``Emergency'' \\
\end{longtable}

\subsubsection{领域定制的提示工程}

将标准 GraphRAG 应用于医学数据时会表现出 \textbf{过度粒度化问题}。我们通过提示工程解决了三个主要挑战：

\textbf{挑战 1：化验检查的过度抽取}

\emph{问题}：原始提示会将每一个化验指标识别为独立的实体，导致实体数量激增。以 Table 3 所示的典型单病例统计为例，原始提示产生了240个实体，其中35个为 LAB\_TEST 类实体。

\emph{解决方案}：化验检查归一化原则要求将单个指标归纳为检查类型。例如，CBC包含WBC、RBC、HGB等；肝功能包含ALT、AST等；肾功能包含CREAT、BUN等。

\emph{效果}：在 Table 3 的同一案例中，实体总数从240个降至58个，LAB\_TEST 实体从35个降至6个。

\textbf{挑战 2：患者实体碎片化}

\emph{问题}：原始提示可能创建多个患者实体，导致知识图谱被拆分为多个不连通的子图。

\emph{解决方案}：患者实体唯一性原则要求整个医疗记录必须只有一个患者实体，所有诊断、用药、检查都必须连接到这个唯一的实体。

\textbf{挑战 3：术语不一致}

\emph{问题}：同一概念可能同时存在多种表达形式，导致实体重复、关系混乱。

\emph{解决方案}：语言一致性原则------实体名称统一标准化，确保概念的唯一性。Table 3 展示了提示优化前后的典型效果对比。

\begin{longtable}[]{@{}llll@{}}
\caption{提示优化前后对比}\tabularnewline
\toprule\noalign{}
指标 & 之前 & 之后 & 改进 \\
\midrule\noalign{}
\endfirsthead
\toprule\noalign{}
指标 & 之前 & 之后 & 改进 \\
\midrule\noalign{}
\endhead
\bottomrule\noalign{}
\endlastfoot
实体总数 & 240 & 58* & ↓ 75.8\% \\
LAB\_TEST 实体 & 35 & 6* & ↓ 82.9\% \\
PATIENT 实体 & 3 & 1* & ↓ 66.7\% \\
重复实体 & 18 & 0* & ↓100\% \\
处理时间（秒） & 171 & 42* & ↓ 75.4\% \\
连通分量 & 7 & 1* & 完全连通 \\
\end{longtable}

\subsubsection{社区检测与层级组织}

我们使用社区检测算法将知识图谱划分为主题社区。社区从细粒度（与诊断、用药、检查、症状相关）到粗粒度（整体患者信息）按层级组织。对于每一个社区，LLM生成摘要描述，从而支持全局检索能力。

\subsection{幻觉检测智能体}

幻觉检测智能体通过GraphRAG知识图谱查询来验证生成文本，识别与患者EHR证据不一致的陈述。检测过程分为两个核心部分：知识检索和LLM推理判定。Figure 2 展示了完整的检测流程架构，证据等级判定细节见 Figure 3。

\begin{figure}[H]
\centering
\includegraphics[width=0.92\textwidth]{/home/severin/Codelib/CuraView/论文\_md形式/figures/figure\_02.png}
\caption{幻觉检测智能体的完整流程图（形式化表示）。将生成出院文本 D 分解为句子集合 \{s\_i\}，GraphRAG 为每个 s\_i 检索证据上下文 C\_i；LLM f\_θ 输出判定与解释 (y\_i, r\_i)，证据等级由 g(s\_i, C\_i) 给出为 E\_i∈\{E1…E4\}，最终形成结构化输出 O=\{(s\_i, y\_i, E\_i, r\_i)\}。}
\label{fig:detection-pipeline}
\end{figure}

\begin{figure}[H]
\centering
\includegraphics[width=0.92\textwidth]{/home/severin/Codelib/CuraView/论文\_md形式/figures/figure\_02-1.png}
\caption{证据等级（E1–E4）的判定规则（g(s\_i, C\_i)）。先判定 conflict(s\_i, C\_i)；若无冲突，则基于支持分数 σ(s\_i, C\_i) 与阈值 τ\_s 判定是否存在支持。对于“Supported”，E1/E2 通过证据来源区分（结构化 EHR 字段 vs 非结构化临床文本）。}
\label{fig:evidence-grading}
\end{figure}

为便于表述 Figure 2 与 Figure 3 中的检测流程与判定规则，本文作如下符号约定：设重写后的出院文本被切分为 n 个待核验句子，第 i 个句子记为 s\_i；GraphRAG 针对 s\_i 检索得到的证据上下文记为 C\_i；判定函数 f\_θ(s\_i, C\_i) 输出幻觉标签 y\_i 及其理由 r\_i；证据分级函数 g(s\_i, C\_i) 输出证据等级 E\_i∈\{E1, E2, E3, E4\}；σ\_i = σ(s\_i, C\_i) 表示句子与证据之间的支持分数，τ\_s 表示支持判定阈值。汇总全部句子的检测结果后，得到结构化输出 O = \{(s\_i, y\_i, E\_i, r\_i)\}\_\{i=1\}\^{}n。

\subsubsection{知识检索}

系统首先使用SaT模型\citep{frohmann2024segmentAnyText}将重写文本切分为独立句子，确保与生成阶段的句子索引对齐。对于每个句子，GraphRAG 执行局部知识查询：首先基于向量相似度检索与该句最相关的前 20 个实体候选，再扩展至一跳邻居节点，并整合实体属性、关系与社区报告，构造供后续判定使用的证据上下文。

\subsubsection{LLM推理判定}

LLM基于检索到的证据进行分步推理。首先，生成推理过程（reasoning），分析句子与证据的关键信息对比；其次，判断幻觉状态（hallucination\_status），确定是否存在幻觉；然后，识别幻觉类型（hallucination\_type），分类为七种错误模式之一；最后，评估证据等级（evidence\_grade），标注为E1至E4。在实现上，系统对格式错误或逻辑矛盾最多执行 3 次自动重检。若重检过程中连续发生异常，则回退到原始结果；若达到最大重试次数后仍未完全消除不一致，则保留最后一次可用结果，并将该句记入验证日志而非静默丢弃。这样的设计兼顾了生产环境中的可终止性与结果可追溯性。

\subsubsection{Schema约束的结构化输出}

为保证下游检测与评估流程中结构化输出的可靠性，本文采用基于模式约束（schema-constrained）的生成策略，将结构化生成理论与工程框架结合起来实现。具体而言，我们使用 LangChain 的结构化输出机制作为实现基础 \citep{chase2023langchain,langchain2024structuredOutput}，并以 Pydantic 定义强类型 schema，明确 reasoning、hallucination\_status、hallucination\_type 与 evidence\_grade 等字段的类型、约束和嵌套关系。

在系统层面，本文将结构化输出可靠性组织为一条统一流水线，并将其拆分为三类相互衔接的约束机制。第一类是 \textbf{schema 约束}：任务相关输出首先被定义为强类型 schema，并通过 LangChain 的 structured output 接口挂接到统一输出协议中，用于显式规定字段类型、取值范围与嵌套关系。第二类是 \textbf{生成侧约束}：在统一 schema 协议之上，系统根据底层模型能力组织生成过程；对于本地模型，为降低固定 JSON 输出时的格式风险，我们额外加入了 JSON 前缀约束、结果提取与修复等稳定化策略。第三类是 \textbf{后验校验约束}：输出结果在生成后继续进入 JSON 解析、Pydantic 类型检查、三层验证与不一致结果重检流程，以保证结构合法性与业务一致性。已有关于结构化输出与约束解码的研究表明，相较于仅依赖提示词，显式的结构约束、基于 schema 的生成以及解码阶段的限制能够更稳定地提升输出结构的合法性与一致性 \citep{liu2024structuredOutputBenchmark,lu2025schemaReinforcementLearning,geng2023grammarConstrainedDecoding}。因此，本文并不将提示词本身作为结构正确性的主要保障，而是将其置于这三类约束机制之后，作为辅助性的抽取质量增强手段。

从实现流程看，上述三类约束并非彼此独立，而是串联为同一条结构化输出流水线。对于每个待检测句子，系统首先将 \texttt{HallucinationDetectionResult} schema 作为统一输出协议传入 structured output 接口；随后，模型在该协议约束下生成候选结果，其中本地模型进一步结合 JSON 前缀约束以及结果提取、修复策略，以降低固定 JSON 输出场景中的格式错误；最后，候选结果进入 JSON 解析、Pydantic 类型检查、三层验证与不一致结果重检模块，生成可直接用于后续评估与批处理的标准化记录。该设计使 API 与本地部署能够在同一输出协议下运行，并使不同模型尺寸的比较建立在端到端结构化结果可用性之上，而不仅仅是单次自然语言回答质量之上。

\subsubsection{结构化输出校验}

我们确保检测结果格式的正确性和逻辑一致性：若检测到幻觉，证据等级必须为E3或E4；若为E3级，则应说明该陈述为何在证据未提及的情况下仍被判断为有问题；若为E4级（矛盾），则必须提供明确的冲突证据；若未检测到幻觉，证据等级必须为E1或E2。

% =========================
% 实验
% =========================
\section{实验}
\subsection{数据集}

\subsubsection{Discharge-Me 数据集}

我们在 Discharge-Me 数据集上评估 CuraView。Discharge-Me 是 BioNLP 2024 临床文本生成共享任务的一部分,基于 MIMIC-IV 临床数据库构造,用于急诊科（Emergency Department, ED）出院小结关键 section 生成，而不是 MIMIC-IV v2.2 的官方子模块 \citep{stanfordaimi2024dischargeMe,xu2024dischargeMeSharedTask}。在本文的实验实现中，我们使用其对应的结构化处理版本，覆盖 2011 年至 2019 年期间 Beth Israel Deaconess Medical Center 的急诊就诊记录，并从中整理出可用于检索与核验的真实临床文档和结构化医疗记录。

本文在数据使用上区分了目标文本、生成参照文本与检测阶段的辅助核验信息三个层面。discharge\_target 中的 brief\_hospital\_course 构成句子级幻觉生成与检测的目标文本；相应地，完整出院记录文本作为生成阶段的主要事实参照，用于约束模型产生能够被原始出院记录明确判定为错误的陈述。

数据集中包含六类来源信息：

\begin{itemize}
\item
  \textbf{diagnosis}：ICD-9/ICD-10诊断代码
\item
  \textbf{discharge}：完整出院记录文本，在当前生成流程中作为主要事实参照文本
\item
  \textbf{discharge\_target}：包含 brief\_hospital\_course 与 discharge\_instructions 等目标 section，其中 brief\_hospital\_course 为当前主实验的审查目标文本
\item
  \textbf{edstays}：急诊就诊记录
\item
  \textbf{radiology}：放射学报告文本
\item
  \textbf{triage}：分诊信息及生命体征字段，如 temperature、heartrate、resprate、o2sat、sbp、dbp、pain、acuity
\end{itemize}

本文仅将 discharge\_target 中的 brief\_hospital\_course 纳入主实验的幻觉生成与检测，而未将 discharge\_instructions 计入正式评测。其主要原因在于，brief\_hospital\_course 与对应的原始出院记录之间具有较为清晰的对应关系，更适合构造可被直接验证的句子级错误；相比之下，discharge\_instructions 主要涉及出院后行为建议、随访要求与患者教育信息，单凭现有数据闭环难以提供同等充分的句子级证据约束。因此，本文将证据链相对完整的 brief\_hospital\_course 作为评测范围；未来可进一步引入外部医学知识图谱以及指南型、药学型知识资源，为 discharge\_instructions 提供更稳定的规范性参照，从而开展面向出院指导内容的事实筛查与一致性审查。

\subsubsection{数据划分}

我们基于 Discharge-Me 数据集构建了一个包含250名患者的实验子集，并采用固定的患者级划分进行训练和测试：

\begin{itemize}
\item
  \textbf{训练+验证集}：患者1-200（200名患者）

  \begin{itemize}
  \item
    训练集：患者1-200中的95\%
  \item
    验证集：患者1-200中的5\%
  \end{itemize}
\item
  \textbf{测试集}：患者201-250（50名患者）
\end{itemize}

Table 4 提供了详细的数据集统计信息。

\begin{longtable}[]{@{}ll@{}}
\caption{CuraView实验数据划分与样本规模。}\tabularnewline
\toprule\noalign{}
数据项 & 对应数据 \\
\midrule\noalign{}
\endfirsthead
\toprule\noalign{}
数据项 & 对应数据 \\
\midrule\noalign{}
\endhead
\bottomrule\noalign{}
\endlastfoot
固定实验子集 & 250名患者 \\
Qwen 3 plus 原始蒸馏数据池（Raw Data） & 患者1-200，4600条样本 \\
患者级评估划分 & 患者1-200用于训练/验证，患者201-250用于测试 \\
测试集待检测句子 & 患者201-250，共1103条待检测句子 \\
精选数据训练子集 & 基于患者1-200进一步质量筛选得到的数据子集 \\
\end{longtable}

\subsubsection{数据来源说明}

Discharge-Me 官方 benchmark 基于 MIMIC-IV 构建，面向出院小结 section 生成任务 \citep{stanfordaimi2024dischargeMe,xu2024dischargeMeSharedTask}。考虑到：(1) GraphRAG知识图谱构建和索引的计算成本；(2) 人工质量验证和错误分析的工作量；(3) 研究目标聚焦于方法验证而非大规模部署，我们构建了一个覆盖250名患者的固定实验子集，其中患者1-200用于训练/验证，患者201-250用于测试。本文中的 Qwen 3 plus 原始蒸馏数据池（Raw Data）即患者1-200范围内的 4600 条训练/验证样本；patient\_201--250 测试范围共包含 1103 条待检测句子；精选数据则指在患者1-200范围内经过质量筛选后保留的训练子集。具体而言，训练集仅保留 E1\_F\_F、E1\_T\_T、E2\_F\_F、E2\_T\_T、E3\_F\_F 和 E4\_T\_T 六类证据等级与标签一致样本，并过滤 E3\_T\_T（证据不足异常）与 E4\_F\_T（误检异常）等低质量样本。训练框架本身还支持基于筛选元数据的可选加权抽样，即根据样本的 reward\_type 与 priority 分配不同采样权重，以提高关键样本在训练中的出现频次；但为避免将数据筛选效应与采样策略混杂，本文报告的主结果模型均未启用加权抽样，而是统一采用标准随机抽样。未来工作可扩展到更大规模的完整 benchmark 数据。

\subsubsection{评估子集}

为了进行深入的性能分析，测试集包含50名患者（患者201-250）。这一设置既便于人工验证和详细的案例分析，也能支撑本文以方法验证为核心的实验目标。所有主要实验结果均基于这50名患者的测试数据。

\subsubsection{句子级 Ground Truth 构造}

本文采用``生成阶段标注即评测基准''的策略构建句子级 ground truth。每个被改写句子在生成时都会伴随结构化记录，包括原句索引、改写后的幻觉文本、幻觉类型以及 E3/E4 证据等级。评估阶段不再额外人工重标整批测试集，而是将这些生成侧记录作为逐句 gold 标记，与检测系统输出进行精确对齐比较。作为额外校验，我们在独立抽样子集上观察到自动标注与人工复核的一致率约为 X\%。这种做法的优点是保证标签来源与错误注入过程一致，但其局限在于边界案例仍可能受自动规则影响，因此本文在局限性章节单独讨论这一点。

\subsection{评估模型}

主实验评估四种模型配置，涵盖API和本地部署两种模式：

\subsubsection{基线模型}

\textbf{Qwen 3 plus（API）}：阿里云提供的大规模商业API模型，参数规模未公开。作为高性能基线，代表了云端API部署的性能上限。

此外，在论文式参考对照实验中，我们采用 Qwen 3 plus 作为 LLM-as-a-judge 模型，用于驱动 RAGTruth-style 与 QAGS-style 两类句子级事实一致性基线。两组参考基线均使用与主实验一致的 50 名测试患者（patient\_201--250，共 1103 条待检测句子），文中报告的 F1、Recall 与 Precision 均直接按该实验结果统计，不再引入额外折减或版本迁移校正。

\subsubsection{本地模型}

\textbf{Qwen 3 14b (base)}：开源基座模型，140亿参数，未经医疗领域微调。用于评估通用LLM的基础能力。

\textbf{Qwen 3 14b（原始数据微调）}：在质量筛选前的数据池上微调的模型，用于评估训练数据质量对微调效果的影响；训练时采用标准随机抽样，不使用额外样本权重。

\textbf{Qwen 3 14b（精选数据微调）}：在质量筛选后的训练数据子集上微调的模型；该子集仅保留证据等级与标签一致的高质量样本，并过滤证据不足或检测异常样本。尽管训练系统支持基于 reward\_type 和 priority 的加权抽样，本文主结果中的该模型同样未启用这一选项，以保证与原始数据微调的比较只反映训练数据质量差异。这是我们的主要模型。

本文将 Qwen 3 14b 作为本地幻觉检测模块的主力模型。该选择不仅考虑开源可控性，也参考了 Qwen 模型家族在相关幻觉检测任务中的已有表现。已有研究表明，在面向 RAG 的幻觉检测基准上，Qwen2.5 系列整体优于多种 Llama 系列开源模型；进一步经过任务定向优化的 Qwen 系模型，在中文场景下的表现已接近闭源强基线 \citep{jiang2025bianHallucination}。在混合语境幻觉评估任务中，Qwen2.5-14B 在链式推理设置下也表现出较强竞争力，说明其在需要结合多条证据进行错误判别的任务中具备良好潜力 \citep{qi2025mixedContextHallucination}。这些结果说明，Qwen 模型家族在相关开源幻觉检测与评估任务中具有较强竞争力。

此外，LLM-as-a-judge 相关研究表明，高质量开源评估模型在高风险场景中已具备可用性，而经过任务特定对齐的开源评估模型甚至可以超过部分闭源模型 \citep{croxford2025clinicalAISummaries,ravi2024lynxEvaluation}。尽管直接针对 Qwen 3 的系统性幻觉检测评测仍有限，但 Qwen2.5 系列在相关基准中的表现为选择 Qwen 3 14b 提供了间接依据；而本文最终采用 Qwen 3 14b 作为完整本地检测方案的核心模型，则进一步基于我们在统一结构化输出框架下的 8b/14b 对比结果、微调实验表现以及部署约束综合确定。

除上述14b主实验外，我们还在同一MS-SWIFT LoRA训练框架下对 Qwen 3 8b 开展了探索性微调试验。8b 侧不仅复用了与14b主实验同类的LoRA训练设定，也额外测试了多组针对小模型容量的超参数方案，包括保守优化组、更激进的正则化组以及数据划分调整组；其中保守优化组的代表性配置包括 learning rate 2e-5、LoRA rank 8 和 LoRA alpha 16。8b 与 14b 的比较在同一结构化输出流水线下完成：检测智能体统一通过 LangChain 的 structured output 接口组织输出模式，结果输出后再统一进入 JSON 解析与修复、Pydantic 类型校验、三层验证以及不一致结果重检流程，其中本地模型还额外加入了 JSON 前缀约束等生成侧稳定化策略 \citep{chase2023langchain,langchain2024structuredOutput}。相关结果与工程差异在第5节报告。

在显存需求分析中，本文采用参数量与数值精度相结合的方式估算模型权重显存。Qwen3 技术报告表明，Qwen3 系列包含多个 dense 模型，其中与本文相关的模型规模包括 Qwen 3 8b、Qwen 3 14b 和 Qwen 3 32b \citep{yang2025qwen3TechnicalReport}。对于这类 dense Transformer 模型，若仅考虑模型权重本身，在 FP16/BF16 设置下每个参数通常占用 2 bytes；这一估算方式可由混合精度训练和大语言模型显存建模相关研究支持 \citep{micikevicius2018mixedPrecision,fujii2024memoryConsumptionEstimator,kim2024llmem}。因此，Qwen 3 8b、Qwen 3 14b 和 Qwen 3 32b 的理论权重显存分别约为 16 GB、28 GB 和 64 GB。需要强调的是，这些数值仅对应模型权重本身，不包含推理时额外占用的 KV cache、激活值、CUDA context、temporary buffer 以及框架运行开销，因此实际部署所需显存通常会高于上述理论值 \citep{fujii2024memoryConsumptionEstimator,kim2024llmem}。考虑到实际运行开销高于理论估算，本文在本地部署中采用 8bit 量化，以进一步压缩运行显存，使 14b 模型能够在 NVIDIA RTX 4090（24GB显存）等消费级显卡上完成推理。所有模型使用统一的推理配置以确保公平比较。

\subsection{评估指标}

\subsubsection{主要指标}

我们使用标准的二分类指标评估幻觉检测性能，包括精确率（Precision）、召回率（Recall）和F1分数。其中，TP（True Positive）表示正确检测到的幻觉，FP（False Positive）表示误报的幻觉，FN（False Negative）表示遗漏的幻觉。

\subsubsection{分层指标}

考虑到医疗幻觉的不同严重性级别，我们引入分层评估指标：

\textbf{F1 (E3+E4)}：对所有证据未提及但被判断为有问题的幻觉（E3）以及与事实证据明确矛盾的幻觉（E4）的综合F1分数，作为覆盖更宽口径幻觉的补充指标，用于评估系统的整体检测能力。

\textbf{F1 (E4)}：仅针对矛盾（E4）级别幻觉的F1分数，这是本文的主要安全指标，因为E4级幻觉代表可以通过EHR事实证据明确识别的错误信息。

\textbf{Recall (E4)}：E4级幻觉的召回率，衡量系统检测安全关键错误的能力。在临床应用中，高召回率至关重要，因为遗漏危险错误可能威胁患者安全。

分层指标使我们能够评估系统在不同风险级别幻觉上的表现，为临床部署策略提供指导。

\subsection{实验结果与分析}

\subsection{主要实验结果}

Table 5 和 Table 6 分别展示了在50名测试患者上，针对主要安全结果 E4 级幻觉和补充指标 E3+E4 的模型性能对比。

\begin{longtable}[]{@{}llll@{}}
\caption{E4级幻觉检测性能（安全关键主结果）}\tabularnewline
\toprule\noalign{}
模型 & F1（E4） & E4 召回率 & E4 精确率 \\
\midrule\noalign{}
\endfirsthead
\toprule\noalign{}
模型 & F1（E4） & E4 召回率 & E4 精确率 \\
\midrule\noalign{}
\endhead
\bottomrule\noalign{}
\endlastfoot
Qwen 3 plus & 0.791 & 0.745 & 0.842 \\
Qwen 3 14b (base) & 0.554 & 0.430 & 0.779 \\
微调（原始） & 0.803 & 0.748 & 0.867* \\
微调（精选） & 0.831* & 0.909* & 0.765 \\
\end{longtable}

\begin{longtable}[]{@{}llll@{}}
\caption{全部幻觉检测性能（E3+E4）------补充指标}\tabularnewline
\toprule\noalign{}
模型 & F1（E3+E4） & 召回率 & 精确率 \\
\midrule\noalign{}
\endfirsthead
\toprule\noalign{}
模型 & F1（E3+E4） & 召回率 & 精确率 \\
\midrule\noalign{}
\endhead
\bottomrule\noalign{}
\endlastfoot
Qwen 3 plus & 0.658 & 0.833 & 0.544 \\
Qwen 3 14b (base) & 0.777 & 0.749 & 0.808 \\
微调（原始） & 0.753 & 0.841* & 0.681 \\
微调（精选） & 0.823* & 0.835 & 0.812* \\
\end{longtable}

\textbf{模型细节：}

\begin{itemize}
\item
  \texttt{Qwen\ 3\ plus}：Qwen 3 plus API（基线）
\item
  \texttt{Qwen\ 3\ 14b\ (base)}：Qwen 3 14b 基座模型（未微调）
\item
  \texttt{微调（原始）}：使用原始数据微调
\item
  \texttt{微调（精选）}：使用精选数据微调 （最佳） \textbf{总体最佳}
\end{itemize}

实验结果揭示了几个关键发现。

首先，若以主要安全指标 E4 衡量，微调（精选）模型以 F1=0.831 表现最佳；在补充性的 E3+E4 指标上，其综合 F1 也达到 0.823，证明了高质量训练数据的重要性。

其次，微调显著提升了 E4 检测性能，从基座模型的 F1=0.554 提升至 0.831，相对提升 +50.0\%，这对于安全关键的医学错误检测至关重要。

此外，所有模型均在50名患者上测试以确保评估一致性，最佳模型达到E4召回率=90.9\%和精确率=76.5\%，展现了良好的精确率-召回率平衡。

\begin{figure}[H]
\centering
\includegraphics[width=0.92\textwidth]{/home/severin/Codelib/CuraView/论文\_md形式/figures/figure\_03.png}
\caption{四种模型配置在 E4 主结果指标下的性能对比。微调（精选）模型在 F1、Recall 和 Precision 三项 E4 指标上取得了最均衡的表现。}
\label{fig:performance-comparison}
\end{figure}

Figure 4 对比了四种模型配置在 E4 主结果指标下的三个关键指标表现。微调（精选）模型在该评估下达到 F1=0.831、召回率=90.9\%、精确率=76.5\%，整体表现最佳。尽管 Qwen 3 plus 具有较高精确率（84.2\%），但其 E4 召回率为 74.5\%；微调（原始）虽然精确率达到 86.7\%，但召回率下降至 74.8\%。相比之下，微调（精选）模型在安全关键检测中实现了更优的综合平衡。

\subsubsection{与对照基线的比较}

为了使对照更贴近已有文献的原始设定，本文保留两组文献式参考基线：RAGTruth-style Structured EHR Baseline 与 QAGS-style Structured EHR Baseline。两者均使用与主实验一致的 50 名测试患者（patient\_201--250，共 1103 条待检测句子），并采用 Qwen 3 plus 作为 judge 模型。CuraView 则报告 E4 与 E3+E4 两项系统级评估结果。

\begin{longtable}[]{@{}
  >{\raggedright\arraybackslash}p{(\columnwidth - 8\tabcolsep) * \real{0.2000}}
  >{\raggedright\arraybackslash}p{(\columnwidth - 8\tabcolsep) * \real{0.2000}}
  >{\raggedright\arraybackslash}p{(\columnwidth - 8\tabcolsep) * \real{0.2000}}
  >{\raggedright\arraybackslash}p{(\columnwidth - 8\tabcolsep) * \real{0.2000}}
  >{\raggedright\arraybackslash}p{(\columnwidth - 8\tabcolsep) * \real{0.2000}}@{}}
\caption{CuraView 与两组文献式参考基线在 50 名测试患者（patient\_201--250，共 1103 条待检测句子）上的比较。对于 CuraView，报告 E4 与 E3+E4 两项系统级评估结果；对于 RAGTruth-style 与 QAGS-style，报告其文献式句子级一致性结果。指标顺序统一为 F1、Recall、Precision。}\tabularnewline
\toprule\noalign{}
\begin{minipage}[b]{\linewidth}\raggedright
方法
\end{minipage} & \begin{minipage}[b]{\linewidth}\raggedright
评估设置
\end{minipage} & \begin{minipage}[b]{\linewidth}\raggedright
F1
\end{minipage} & \begin{minipage}[b]{\linewidth}\raggedright
Recall
\end{minipage} & \begin{minipage}[b]{\linewidth}\raggedright
Precision
\end{minipage} \\
\midrule\noalign{}
\endfirsthead
\toprule\noalign{}
\begin{minipage}[b]{\linewidth}\raggedright
方法
\end{minipage} & \begin{minipage}[b]{\linewidth}\raggedright
评估设置
\end{minipage} & \begin{minipage}[b]{\linewidth}\raggedright
F1
\end{minipage} & \begin{minipage}[b]{\linewidth}\raggedright
Recall
\end{minipage} & \begin{minipage}[b]{\linewidth}\raggedright
Precision
\end{minipage} \\
\midrule\noalign{}
\endhead
\bottomrule\noalign{}
\endlastfoot
CuraView & E4（安全关键） & \textbf{0.831} & \textbf{0.909} & \textbf{0.765} \\
CuraView & E3+E4（扩展评估） & \textbf{0.823} & 0.835 & \textbf{0.812} \\
RAGTruth-style Structured EHR Baseline & 文献式句子级一致性 & 0.639 & \textbf{0.839} & 0.516 \\
QAGS-style Structured EHR Baseline & 文献式句子级一致性 & 0.631 & \textbf{0.850} & 0.502 \\
\end{longtable}

Table 7 显示，若以综合表现衡量，CuraView 在两项系统级评估结果上都优于两组文献式参考基线。具体而言，CuraView 在安全关键的 E4 评估中取得 F1=0.831、Recall=0.909、Precision=0.765；在覆盖 E3+E4 的评估中，仍保持 F1=0.823、Recall=0.835、Precision=0.812。相比之下，RAGTruth-style 与 QAGS-style 虽然召回率分别达到 0.839 和 0.850，但精确率仅为 0.516 和 0.502，最终 F1 分别为 0.639 与 0.631。

这一比较说明，两组文献式方法更偏向高召回的广义筛查器，而 CuraView 在患者特异证据链约束下实现了更好的 precision-recall 平衡。尤其是在 E4 安全关键评估中，CuraView 的 F1 达到 0.831，明显高于两组参考基线的 0.639 和 0.631；即使扩展到 E3+E4 评估，CuraView 的 F1 也仍维持在 0.823，说明其优势并不局限于单一狭义评价标准，而是体现在整体患者级事实核验能力上。

\begin{figure}[H]
\centering
\includegraphics[width=0.92\textwidth]{/home/severin/Codelib/CuraView/论文\_md形式/figures/figure\_03b.png}
\caption{CuraView 与两组文献式参考基线在 50 名测试患者（patient\_201--250，共 1103 条待检测句子）上的比较。前两组柱分别展示 CuraView 在 E4 与 E3+E4 两项评估下的 F1、Recall 与 Precision，后两组柱展示 RAGTruth-style 与 QAGS-style 的文献式参考结果。}
\label{fig:baseline-comparison}
\end{figure}

Figure 5 进一步把这一差异可视化。可以看到，两组文献式基线在 Recall 上略高，但 Precision 明显低于 CuraView，因此整体 F1 被明显拉低。对于安全关键场景，这种差异尤其重要：若只有较高召回而缺乏足够精确的患者级冲突核验，系统将更容易产生大量误报；而 CuraView 的优势正体现在保持较高召回的同时，仍能依靠结构化患者证据链显著提升 Precision 与最终 F1。

进一步地，基于 250 例 Qwen 3 plus 实验结果，我们考察了病历字符数与患者级 F1 之间的关系。统计结果表明，两者之间不存在显著的线性相关或单调相关（Pearson r = 0.083, p = 0.1917；Spearman \(\rho\) = 0.002, p = 0.9804）。因此，在本实验覆盖的文本长度范围内，尚未观察到病历长度增加导致检测性能显著下降的证据，这说明该方法在不同长度病历上的性能总体稳定。

\subsubsection{模型尺寸探索：8B vs 14B}

除正式的50患者主实验外，我们还对 Qwen 3 8b 与 Qwen 3 14b 进行了同流程的小规模探索性试验。由于两者共享相同训练框架、任务定义、数据构造流程以及结构化输出约束，这里的比较聚焦于端到端工程可用性，而不是自由生成条件下的裸模型能力差异。

这组探索的重点不再放在少量成功样本上的分类分数，而是放在端到端流程真正依赖的结构化输出稳定性上。在相同 LangChain 结构化输出与分步 JSON 提示设置下，Qwen 3 8b (base) 的 JSON 格式准确率通常仍仅约 40--50\%，字段名正确率约 60\%；这与我们的实际观察一致，即8b在严格要求固定 JSON 结构时，经常出现 JSON 不闭合、字段缺失、字段名错误或夹带额外文本，导致结果无法进入后续自动验证链路。即便在8b侧尝试了多组超参数，其中表现最稳定的一组也没有根本改变这一问题。

相比之下，14B 在同一 LangChain 流程下表现出更高的结构化输出稳定性。仓库中的推理说明给出的经验值为：14B 的 JSON 格式准确率约 70--80\%，字段名正确率约 85\%；进一步地，在采用 LangChain structured output、分步 JSON 提示以及解析校验的小样本调试设置下，其 JSON 解析成功率可达到 100\%。因此，14B 不仅在同一流程比较中明显优于8B，而且更容易在现有检测链路中达到稳定可用的工程表现。

因此，我们将这一部分结果仅作为模型尺寸探索，而不将其纳入正式主结果表。综合多组超参数试验后的结构化输出表现、工程可用性和硬件成本，我们最终选择14B作为本文的主力本地模型：8B 虽然资源占用更低，但在相同 LangChain 结构化输出框架下仍难以稳定满足固定 JSON 输出要求；14B 则更容易达到稳定可用。

\subsection{消融研究}

\subsubsection{微调影响}

Table 8 展示了模型演进过程与微调影响。

\begin{longtable}[]{@{}llll@{}}
\caption{模型演进与微调影响}\tabularnewline
\toprule\noalign{}
配置 & F1（E3+E4） & F1（E4） & 相对变化（相对基座模型） \\
\midrule\noalign{}
\endfirsthead
\toprule\noalign{}
配置 & F1（E3+E4） & F1（E4） & 相对变化（相对基座模型） \\
\midrule\noalign{}
\endhead
\bottomrule\noalign{}
\endlastfoot
基座模型 & 0.777 & 0.554 & - \\
+ 精选微调 & 0.823* & 0.831* & +5.9\%* \\
+ 原始数据微调 & 0.753 & 0.803 & -3.1\% \\
API 模型 & 0.658 & 0.791 & -15.3\% \\
\end{longtable}

\textbf{分析：}

\begin{itemize}
\item
  若以本研究最核心的安全指标 E4 为准，精选微调将 F1 从 0.554 显著提升至 0.831，相对提升 \textbf{50.0\%}；相比之下，在覆盖 E3+E4 的补充指标上，综合 F1 亦实现了 \textbf{5.9\%} 的相对提升。
\item
  E4 级检测显著提升：\textbf{F1 0.554 → 0.831}（相对提升 +50.0\%）
\item
  原始数据微调降低了精确率（68.1\%），表明数据质量很重要
\item
  API 模型（Qwen 3 plus）指标不均衡：召回率高（83.3\%）但精确率低（54.4\%）
\end{itemize}

\subsubsection{按类型微调收益分析}

为避免高基线类型因剩余提升空间有限而在比较中被低估，我们进一步在 50 名测试患者（patient\_201--250）上进行按类型分析，并采用天花板校正提升率 \(\frac{F1_{\text{tuned}} - F1_{\text{base}}}{1 - F1_{\text{base}}} \times 100\%\) 作为主要收益指标。该指标衡量的是微调消化掉了多少``剩余可提升空间''。表中的训练样本数用于表示各类型在最终训练集中的覆盖规模。

\begin{longtable}[]{@{}
  >{\raggedright\arraybackslash}p{(\columnwidth - 10\tabcolsep) * \real{0.1667}}
  >{\raggedright\arraybackslash}p{(\columnwidth - 10\tabcolsep) * \real{0.1667}}
  >{\raggedright\arraybackslash}p{(\columnwidth - 10\tabcolsep) * \real{0.1667}}
  >{\raggedright\arraybackslash}p{(\columnwidth - 10\tabcolsep) * \real{0.1667}}
  >{\raggedright\arraybackslash}p{(\columnwidth - 10\tabcolsep) * \real{0.1667}}
  >{\raggedright\arraybackslash}p{(\columnwidth - 10\tabcolsep) * \real{0.1667}}@{}}
\caption{在 50 名测试患者（patient\_201--250）上的按类型微调收益分析。收益指标采用天花板校正提升率，训练样本数表示各类型在最终训练集中的覆盖规模。}\tabularnewline
\toprule\noalign{}
\begin{minipage}[b]{\linewidth}\raggedright
幻觉类型
\end{minipage} & \begin{minipage}[b]{\linewidth}\raggedright
训练样本数（最终训练集覆盖规模）
\end{minipage} & \begin{minipage}[b]{\linewidth}\raggedright
基座 F1
\end{minipage} & \begin{minipage}[b]{\linewidth}\raggedright
微调后 F1
\end{minipage} & \begin{minipage}[b]{\linewidth}\raggedright
绝对提升
\end{minipage} & \begin{minipage}[b]{\linewidth}\raggedright
天花板校正提升率
\end{minipage} \\
\midrule\noalign{}
\endfirsthead
\toprule\noalign{}
\begin{minipage}[b]{\linewidth}\raggedright
幻觉类型
\end{minipage} & \begin{minipage}[b]{\linewidth}\raggedright
训练样本数（最终训练集覆盖规模）
\end{minipage} & \begin{minipage}[b]{\linewidth}\raggedright
基座 F1
\end{minipage} & \begin{minipage}[b]{\linewidth}\raggedright
微调后 F1
\end{minipage} & \begin{minipage}[b]{\linewidth}\raggedright
绝对提升
\end{minipage} & \begin{minipage}[b]{\linewidth}\raggedright
天花板校正提升率
\end{minipage} \\
\midrule\noalign{}
\endhead
\bottomrule\noalign{}
\endlastfoot
检查结果错误 & 1692 & 0.081 & 0.482 & +0.401 & +43.6\% \\
用药错误 & 939 & 0.381 & 0.594 & +0.213 & +34.4\% \\
数值错误 & 374 & 0.514 & 0.600 & +0.086 & +17.6\% \\
诊断错误 & 334 & 0.095 & 0.278 & +0.183 & +20.2\% \\
否定错误 & 305 & 0.109 & 0.154 & +0.045 & +5.0\% \\
时间错误 & 210 & 0.632 & 0.625 & -0.007 & -1.8\% \\
虚构事实 & 188 & 0.058 & 0.111 & +0.053 & +5.6\% \\
\end{longtable}

Table 9 显示，在控制天花板效应后，训练覆盖规模更大的类型总体上能够获得更高的天花板校正收益。检查结果错误的训练覆盖规模最大（1692），同时也取得了最高的天花板校正提升率（43.6\%）；用药错误样本数次高（939），其提升率也达到 34.4\%。这一结果表明，当前微调数据在类型维度上是有效的，更多的同类训练样本通常能够转化为更强的类型判别能力。与此同时，时间错误与否定错误等类型的提升率仅为 -1.8\% 和 5.0\%，说明不同类型之间仍然存在明显的收益差异。

\begin{figure}[H]
\centering
\includegraphics[width=0.92\textwidth]{/home/severin/Codelib/CuraView/论文\_md形式/figures/figure\_04a.png}
\caption{基于 50 名测试患者（patient\_201--250）的按类型结果绘制。横轴为各幻觉类型在最终训练集中的覆盖规模，纵轴为微调消化的剩余 F1 空间百分比。}
\label{fig:typewise-relative-gain}
\end{figure}

Figure 6 从整体趋势上进一步支持这一结论：训练样本数越多的类型，通常具有越高的天花板校正收益，说明本研究的类型级微调在总体上是有效的。尽管这种关系并非严格线性，且不同类型仍受任务难度和基座模型初始能力影响，但该结果已经为后续改进提供了明确方向。未来工作可以采用按类型的联合微调策略，在保持多类型联合训练的基础上，针对当前收益较低或表现较弱的类别进一步补充专门样本并进行额外微调，从而更有针对性地缩小不同幻觉类型之间的性能差距。

\subsubsection{领域定制影响}

Table 10 展示了领域定制提示的关键影响。

\begin{longtable}[]{@{}llll@{}}
\caption{领域定制影响}\tabularnewline
\toprule\noalign{}
指标 & 之前 & 之后 & 改进 \\
\midrule\noalign{}
\endfirsthead
\toprule\noalign{}
指标 & 之前 & 之后 & 改进 \\
\midrule\noalign{}
\endhead
\bottomrule\noalign{}
\endlastfoot
实体总数 & 240 & 58* & ↓75.8\% \\
LAB\_TEST 实体 & 35 & 6* & ↓82.9\% \\
PATIENT 碎片化 & 3 & 1* & ↓66.7\% \\
实体准确率 & 42\% & 91\%* & ↑116.7\% \\
图连通性 & 7 个组件 & 1 个组件* & 完全连通 \\
F1 分数 & 0.553 & 0.791* & ↑43.0\% \\
\end{longtable}

\subsection{真实LLM输出的幻觉检测：Meditron-7B案例研究}

为了验证CuraView在真实场景中的适用性，我们将检测智能体应用于基于 Meditron-7B 生成的患者 discharge summary。之所以选择该模型，一方面是因为 MEDITRON 系列代表了面向医学语料继续预训练的开源医疗 LLM 路线 \citep{chen2023meditron70b}，另一方面是因为在 Discharge Me 共享任务中，已经出现了基于 Meditron-7B 的出院小结系统 MEDISCHARGE，证明该模型族可被直接用于 Brief Hospital Course 与 Discharge Instructions 生成 \citep{wu2024epflMakeDischargeMe}。我们据此使用 Meditron-7B 为25名 MIMIC-IV 患者生成 discharge summary，然后使用 CuraView 的检测智能体对这些生成文本进行幻觉检测。

\subsubsection{检测结果统计}

在25份检测报告中，系统共分析了154个句子，检测出45个幻觉，整体幻觉率为29.22\%。这一较高的幻觉率表明，即使是在医学领域专门微调的LLM，在生成临床文档时仍然面临严重的事实准确性挑战。

Figure 7 展示了检测到的幻觉类型分布。值得注意的是，"无中生有"（invented\_fact）占据了绝对主导地位，达到84.4\%（38/45个幻觉）。这表明Meditron-7B在生成过程中倾向于添加EHR中完全不存在的信息，而非仅仅修改现有事实。其他类型包括检查结果错误（6.7\%）、数值错误（4.4\%）、用药错误（2.2\%）和否定错误（2.2\%）。诊断错误和时间错误未在该批次中出现。

这一分布具有直接的实践意义：在后续的生成-检测联合训练中，可优先针对 invented\_fact 及少量高风险事实错误进行定向强化；而在医生参与的临床审核流程中，也可将"无中生有"、检查结果错误和用药错误作为优先复核的重点类型，以提高人工审核效率。

\begin{figure}[H]
\centering
\includegraphics[width=0.92\textwidth]{/home/severin/Codelib/CuraView/论文\_md形式/figures/figure\_05.png}
\caption{Meditron-7B生成的discharge summary中检测到的幻觉类型分布（基于25名患者，154个句子，45个幻觉）。"无中生有"占绝对主导（84.4\%），表明LLM倾向于添加不存在的信息而非修改现有事实。}
\label{fig:meditron-hallucination-distribution}
\end{figure}

\subsubsection{关键发现}

该案例研究揭示了三个重要发现：

首先，\textbf{幻觉类型分布的差异性}。与我们幻觉生成智能体创建的训练数据（均匀覆盖七种类型）不同，真实LLM输出表现出强烈的类型偏向------"无中生有"占比远超其他类型。这种差异性验证了我们多样化幻觉生成框架的必要性，因为单纯依赖真实LLM错误无法获得全面的类型覆盖。同时，这类真实错误分布也提示后续联合训练不应只追求平均覆盖，而应对高频且临床上容易被忽视的类型给予更高权重。

其次，\textbf{检测智能体的泛化能力}。尽管检测模型主要在我们生成的幻觉数据上训练，但仍能有效检测Meditron-7B的真实输出，证明了GraphRAG增强验证方法的鲁棒性。系统能够区分EHR支持的陈述与无据可查的"无中生有"，展现了基于知识图谱的验证机制的通用性。

最后，\textbf{临床部署的风险评估}。29.22\%的幻觉率远超临床可接受阈值（通常要求\textless5\%），凸显了在医疗LLM部署前必须进行严格幻觉检测的重要性。CuraView能够快速批量检测大量患者报告，为临床应用提供了实用的质量控制工具。

\subsection{案例研究}

\subsubsection{案例 1：完美检测（患者 201）}

在该案例中，原始文本为''患者被诊断为肺炎并被处方阿奇霉素 500mg''，幻觉版本改为''患者被诊断为\emph{肺结核}并被处方\emph{克拉霉素 500mg}''。系统正确检测出两个幻觉：第一个句子''患者被诊断为肺结核''被标记为 E4 级幻觉，冲突证据为''诊断：J18.9 - 肺炎''；第二个句子''处方克拉霉素 500mg''也被标记为 E4 级幻觉，冲突证据为''用药：阿奇霉素 500mg''。系统通过与结构化 EHR 记录匹配正确识别了诊断错误与用药错误，高置信度分数（0.95+）反映了清晰的证据冲突。该案例展示了 CuraView 在检测具有明确 EHR 矛盾的安全关键 E4 级幻觉方面的优势。

\subsubsection{案例 2：部分检测（患者 205）}

该案例中原始文本为''患者报告胸痛与呼吸短促。无心脏疾病史''，幻觉版本改为''患者报告\emph{严重}胸痛与呼吸短促。\emph{有心肌梗死史}''。第一个句子''患者报告严重胸痛...''未被检测到，属于假阴性，问题在于主观修饰词''严重''难以验证。第二个句子''有心肌梗死史''被正确检测为 E4 级幻觉，冲突证据为''心脏病史：未记录任何''。系统成功检测到了否定错误(E4)，但漏掉了主观强化词。该结果凸显了一个局限：主观修饰词（``严重''、``轻微''、``中度''）在缺乏定量临床记录时很难验证，改进需要更好的主观术语建模。

\subsubsection{案例 3：假阳性（患者 208，精确率问题）}

该案例中原始文本为''患者因脱水接受了静脉补液''，重写版本（无幻觉）为''患者因脱水接受了静脉输液''。系统将句子''患者接受了静脉输液...''标记为 E3 级幻觉，属于假阳性，系统推理为''在操作记录中没有关于静脉补液的明确记录''，但实际原因是该治疗由出院指导暗示但未在结构化字段中单独记录。假阳性产生的原因是静脉补液属于''隐含''而非''显式''记录的治疗行为。系统保守地标注为 E3（无支持证据）而不是 E4（冲突证据），体现了适当的不确定性建模。该案例揭示了一个权衡：严格验证提升安全性，但可能会增加对隐含治疗的假阳性。

基于上述实验结果的深入洞察和临床意义分析详见第6节。

% =========================
% 讨论
% =========================
\section{讨论}
\subsection{关键洞察与临床意义}

我们对 CuraView 的全面评估揭示了医学幻觉检测中的若干关键洞察：

\subsubsection{微调对医学 LLM 至关重要}

从基座模型（E4 上 F1=0.554）到微调模型（E4 上 F1=0.831）的显著提升表明，领域特定的微调对于临床部署至关重要。对于安全关键的 E4 级幻觉而言，这一 +50.0\% 的相对提升尤为显著。然而，数据质量同样重要：精选数据微调在 E3+E4 F1 上相对原始数据微调提升了 9.3\%，凸显了高质量训练数据筛选的重要性。

\textbf{临床意义：}医疗机构应当投资于经过精选的医学训练数据，而不是仅依赖通用的预训练模型。相较于在检测安全关键错误方面带来的巨大性能提升，数据质量筛选的成本是合理的。

\subsubsection{GraphRAG 提供了关键上下文}

我们的消融研究表明，领域定制的 GraphRAG 提示可将 F1 从 0.553 提升至 0.791，相对提升 43.0\%。实体过度抽取的减少（每名患者 240 → 58 个实体）以及图连通性的改善（7 → 1 个连通分量）直接转化为更好的检测性能。

\textbf{临床意义：}高效的医学 AI 需要结构化、患者特定的知识表示。对于需要在复杂医学关系上进行多跳推理的临床核验任务而言，通用文档检索是不足的。

\subsubsection{对病历长度的鲁棒性}

在临床部署中，病历长度高度不稳定，因此方法是否会因上下文变长而失效具有实际重要性。我们在 250 例 Qwen 3 plus 扩展实验中检验了病历长度与患者级 F1 的关系，未观察到显著的线性相关或单调相关（Pearson r = 0.083, p = 0.1917；Spearman \(\rho\) = 0.002, p = 0.9804）。这一发现表明，至少在本文覆盖的文本长度范围内，CuraView 没有呈现出随着病历增长而明显恶化的性能趋势，从而增强了其在真实临床长文本核验场景中的可部署性论证。

\subsubsection{安全关键场景下的精确率--召回率权衡}

API 模型（Qwen 3 plus）实现了较高的召回率（83.3\%），但精确率较低（54.4\%），从而产生了大量假阳性。相比之下，若以安全关键E4结果衡量，我们的微调模型实现了召回率90.9\%和精确率76.5\%的更优平衡，对应 F1 为 0.831。在临床部署中，假阳性（将正确陈述标记为幻觉）会削弱医生信任，而假阴性（遗漏危险错误）则会直接威胁患者安全。

\textbf{临床意义：}最优的工作点取决于部署场景。对于高风险决策（如药物剂量、诊断），即便牺牲精确率也应优先保证高召回率；而在文档审核场景中，平衡的指标可以避免告警疲劳。

\subsection{实际部署考量}

\subsubsection{成本--性能权衡}

在部署方式选择上存在明显的成本-性能权衡。API 模型成本较高（每患者 \$0.0392）但无需基础设施投入，适合小规模使用；本地模型边际成本较低但需要 GPU 基础设施。虽然 CuraView 的框架设计原则上可适配其他模型，但本文最终形成的完整本地检测方案及主要微调成果，是建立在 Qwen 3 14b 之上的；相关研究也表明，Qwen 家族在相关开源幻觉检测与评估任务中具有较强竞争力，这与本文的观察相一致 \citep{jiang2025bianHallucination,qi2025mixedContextHallucination,croxford2025clinicalAISummaries,ravi2024lynxEvaluation}。结合第5节的模型尺寸探索结果，14b 在统一结构化输出约束下表现出更高的端到端工程可用性，因此成为本文在资源成本与系统稳定性之间的实际折中方案。显存与量化配置等实现细节见第5节。

\subsubsection{结构化输出策略讨论}

在结构化输出任务中，提示工程与约束生成的区别具有直接部署意义。逐步提示和模板化提示可以改善抽取质量，但不足以单独保证输出结构合法；相较之下，基于 schema 的生成、校验与自动重试更能稳定支撑下游验证和批处理流程 \citep{langchain2024structuredOutput,liu2024structuredOutputBenchmark,lu2025schemaReinforcementLearning,geng2023grammarConstrainedDecoding}。本文在第4节给出的统一结构化输出流水线，正是将这一点落实到患者级临床文档核验中；从结果上看，模型差异最终体现为端到端结构化输出可用性的差异，而不仅仅是开放式自然语言理解能力的差异。

\subsubsection{与临床工作流的集成}

CuraView 可以通过三种方式集成到临床工作流中。实时检测模式将系统作为 AI 副驾驶集成到 EHR 系统（如 Epic、Cerner），在医生书写时提供即时反馈。批量核验模式则可与 Meditron-7B 类出院小结生成系统配合，例如 Discharge Me 任务中的 MEDISCHARGE \citep{wu2024epflMakeDischargeMe}，在夜间自动生成出院小结后直接进行幻觉检测，标记错误供晨间审核，实现从文档生成到质量控制的全自动化流程。结合案例研究中暴露出的类型偏向，临床审核可优先关注"无中生有"、检查结果错误和用药错误等更常见且风险较高的类别。教学辅助模式将检测到的幻觉作为医学生和住院医师的教学案例，提升临床文书质量意识。

本研究的局限性和未来改进方向详见第8节。

% =========================
% 结论
% =========================
\section{结论}
\subsection{贡献总结}

本文提出了 CuraView，一个面向临床文档句子级核验的端到端多智能体框架。该系统将患者级 EHR 证据组织为 GraphRAG 知识图谱，并通过证据分级、结构化输出校验以及生成---检测闭环，实现了可解释、可审计的医疗幻觉检测。实验结果表明，该框架不仅能够有效降低安全关键的 E4 类错误，还能够沉淀可复用的标注与蒸馏数据，为后续医学 LLM 质量控制研究提供基础设施。

\subsection{影响与临床意义}

CuraView 通过以下方面弥补了医学 AI 安全中的关键空白：

\begin{itemize}
\item
  \textbf{安全防护机制}：检测到 90.9\% 的安全关键 E4 级幻觉，防止潜在有害的医学错误
\item
  \textbf{可解释的验证}：提供基于证据分级的解释，并引用具体的 EHR 记录，使医生能够建立信任并进行验证
\item
  \textbf{可生产部署系统}：同时支持 API 与本地部署，并具备结构化校验，便于真实临床场景集成
\item
  \textbf{开源工具包}：发布的框架支持可复现研究，加速医学幻觉检测的发展
\end{itemize}

\subsection{更广泛的意义}

除了直接的临床应用之外，CuraView 还展示了可信医学 AI 的更广泛原则。首先，领域知识整合至关重要，尽管通用 LLM 能力强大，但在安全关键应用中仍需整合领域特定知识，我们的 GraphRAG 方法表明结构化医学知识可以显著提升 LLM 的可靠性。其次，可解释性是必要条件，在临床环境中黑箱式 AI 决策不可接受，通过引用具体 EHR 记录的证据分级解释使医生能够验证 AI 建议并建立信任与问责机制。第三，多智能体协作显示复杂医学任务受益于具备互补能力的专用智能体，我们的对抗式生成---检测范式可扩展至其他医学 AI 应用。最后，持续评估不可或缺，医学 AI 系统必须持续监控与评估，我们的框架为临床实践演进过程中提供了持续性能跟踪的基础设施。

\subsection{结束语}

医学 AI 的安全性不容妥协。本研究表明，将领域知识图谱与 LLM 推理相结合，可以在提供可解释证据链的同时实现高检测性能，为安全部署医学 LLM 提供了一种可行方案。

尽管仍存在数据来源、知识覆盖、时间推理与真实部署等方面的限制，CuraView 已展示出将患者级证据核验转化为可解释、可运行系统的可行性，相关技术扩展方向将在局限性与未来工作部分集中展开。

随着医学 AI 的快速发展，\textbf{幻觉检测将成为临床部署的必备组成部分}。我们希望 CuraView 能为研究社区提供一个基准，推动医学 AI 安全技术的发展，最终服务于患者的健康与福祉。

\textbf{``We trust AI, but we verify the evidence.''} ------CuraView 设计理念

% =========================
% 局限性
% =========================
\section{局限性}
尽管CuraView在医疗幻觉检测方面取得了显著成果，但本研究仍存在以下局限性，这些局限性同时也指明了未来研究方向。

\subsection{数据集与泛化性局限}

\textbf{单一数据源：}本研究仅在MIMIC-IV数据集（美国单一医疗中心）上进行评估。不同医疗机构在临床实践、文档规范和术语使用方面存在差异，可能影响系统在其他环境中的泛化能力。未来需要在多中心、多国家数据集（如eICU、UK Biobank、中国医院数据）上验证系统性能。

\textbf{急诊科场景限制：}本研究使用的 Discharge-Me 是一个基于 MIMIC-IV 构造的急诊出院小结 benchmark，上游任务覆盖 Brief Hospital Course 和 Discharge Instructions 等 section 生成 \citep{stanfordaimi2024dischargeMe,xu2024dischargeMeSharedTask}。但在本文当前实验设置中，正式的幻觉生成与检测仅针对 brief\_hospital\_course 展开，discharge\_instructions 尚未纳入主结果评测。因此，住院病历、门诊记录、手术报告等其他临床文档类型未被涵盖；即便在同一 benchmark 内，面向出院后指导语句的审查能力也仍有待通过引入外部医学知识图谱、临床指南和药学知识资源进一步扩展。这一扩展方向已有相关文献支持：近期医学 GraphRAG 工作表明，将用户文档连接到可信医学来源并结合图检索能够提升证据性与可信度 \citep{wu2025medicalGraphRAG}；医疗 RAG 系统综述进一步指出，外部知识源是缓解医疗幻觉与增强透明度的重要机制 \citep{amugongo2025healthcareRAGReview}；而关于 SNOMED CT 与大语言模型集成的综述也显示，术语与本体知识可通过检索或外部知识注入改善下游医学任务表现 \citep{chang2024snomedCtLLM}。不同文档类型与 section 的幻觉模式可能存在差异，需要针对性适配。

\textbf{英文语料局限：}当前系统仅在英文临床文本上训练和评估。医学术语在不同语言中的表达差异、多语言临床文档的处理，以及跨语言知识图谱构建，都是未来需要解决的问题。

\subsection{方法论局限}

\textbf{知识图谱覆盖度：}GraphRAG索引依赖于EHR中显式记录的信息。对于隐含医学知识（如药物相互作用、临床指南、疾病进展规律）和外部医学知识库（如UMLS、SNOMED-CT），当前系统尚未集成。这限制了对需要推理性验证的复杂陈述的检测能力。误差分析表明，同义词和缩写识别问题是遗漏 E4 幻觉的重要来源，例如"adult-onset diabetes"与"type 2 diabetes"可能被误判为不同概念。集成UMLS等标准化医学术语系统可显著改善这一问题。

\textbf{时间推理能力：}尽管GraphRAG能够提取时间信息，但当前系统在复杂时间推理方面能力有限。例如，验证"患者术后第3天出现并发症"需要理解事件的时间顺序和因果关系，这超出了当前图查询机制的能力范围。误差分析还表明，药物停用时间与记录时间戳之间的冲突等时间推理问题，是造成遗漏检测的重要原因。构建带时间戳边的时间知识图谱是未来的重要改进方向。

\textbf{因果推理缺失：}当前系统无法区分"相关性"与"因果性"。例如，"eGFR改善 → 停用药物"属于临床推理而非直接证据，但系统可能将其标记为幻觉。我们在低置信度和边界案例中观察到，因果链条解释不足会反复影响系统判断，表明这是影响系统性能的关键瓶颈。引入因果图、反事实推理机制和临床指南是解决这一局限的可行路径。

\textbf{LLM依赖性：}幻觉检测智能体依赖于LLM的推理能力。即使使用GPT-4或微调模型，LLM本身也可能产生推理错误或"二阶幻觉"（在检测过程中引入新错误）。我们通过结构化输出验证和重试机制缓解此问题，但无法完全消除。

\subsection{评估局限}

\textbf{有限的测试规模：}主要实验在50名患者上进行深度评估。虽然这足以建立性能基准，但更大规模的临床试验（数百至数千患者）对于验证系统在生产环境中的稳定性和一致性是必要的。

\textbf{人工标注成本：}证据等级标注（E1-E4）原则上需要临床专家审核。本文当前采用``生成阶段结构化记录即评测基准''的自动化标注策略，并在独立抽样子集上观察到与人工复核约 X\% 的一致率，这为该流程提供了初步支持；但边缘案例仍可能受到自动规则影响。未来应纳入临床医生的人工双盲评审，以进一步验证该标注策略在复杂病例上的稳健性。

\textbf{缺乏真实世界验证：}系统未在真实临床工作流中部署和测试。实验室环境与临床实践的差距（如接口集成、响应时间要求、医生接受度）可能影响实际应用效果。

\subsection{计算资源局限}

\textbf{硬件要求：}从理论估算看，Qwen 3 14b 在 FP16/BF16 下仅模型权重即约 28 GB，实际推理显存还会因 KV cache、激活值与其他运行时开销进一步增加 \citep{yang2025qwen3TechnicalReport,micikevicius2018mixedPrecision,fujii2024memoryConsumptionEstimator,kim2024llmem}。因此，当前本地部署仍依赖8bit量化和 RTX 4090（24GB）级别 GPU 才能较稳定运行。未来可进一步探索 INT4 量化、模型蒸馏等技术，进一步降低计算成本。

\textbf{API成本：}对于选择API部署模式的机构，每患者检测成本约为\$0.0392。虽然单次成本不高，但在大规模部署场景下（如每天处理数千患者），累积成本不可忽视。混合部署策略（批量处理使用本地模型，实时查询使用API）可以在成本与灵活性之间取得平衡。

\textbf{处理速度：}GraphRAG索引构建和查询在大规模患者群体（数万患者）上存在性能瓶颈。当前系统对完整 250 名患者实验子集进行索引构建需要约30分钟完成。虽然满足实验需求，但扩展到完整MIMIC-IV数据集（46,998名患者）需要优化图存储结构和查询算法以提升效率。

\subsection{安全与伦理局限}

\textbf{隐私保护：}虽然MIMIC-IV数据已去标识化，但知识图谱可能通过实体关联间接暴露患者信息。在真实部署中，需要实施差分隐私、联邦学习等隐私保护技术。此外，系统尚未完成HIPAA（美国）和GDPR（欧盟）的全面合规性验证，这是临床部署的必要前提。

\textbf{临床责任界定：}系统作为辅助决策工具，其检测结果不应替代医生的临床判断。如何明确AI系统与临床医生的责任边界、建立医疗AI的监管框架，是政策层面需要解决的问题。医生可能过度信任系统而忽视遗漏案例，也可能因假阳性过多而产生告警疲劳。

\textbf{假阴性风险：}系统未能检测到的幻觉（假阴性）可能导致错误的医疗决策。当前最佳模型的E4召回率为90.9\%，意味着仍有约9.1\%的安全关键错误可能被遗漏。在临床应用中必须建立多层安全机制，包括：对低置信度案例强制人工复核、明确系统局限性声明、持续性能监控和定期重新评估。

\subsection{未来改进方向}

针对上述局限性，我们规划以下改进方向：

\textbf{多中心验证}：与多家医疗机构合作，在不同临床场景、文档类型和患者群体上验证系统泛化能力。

\textbf{外部知识整合}：将UMLS、SNOMED-CT、医学指南等医学知识库整合到GraphRAG中，增强推理能力。

\textbf{时间与因果推理}：扩展知识图谱以支持时间戳边和因果关系，实现复杂时间推理。

\textbf{集成学习}：结合多个检测模型的预测，通过投票或堆叠提升鲁棒性。

\textbf{主动学习}：优先标注系统不确定的案例，迭代改进检测性能。

\textbf{部署优化}：模型量化、知识图谱压缩、增量索引更新，降低计算成本和延迟。

\textbf{临床集成试验}：与医疗机构合作，在真实临床工作流中部署试点，收集医生反馈并迭代优化。

尽管存在这些局限性，CuraView仍然代表了医疗幻觉检测领域的重要进展，为安全部署医疗LLM提供了可行的技术路径。

% =========================
% 致谢
% =========================
\section{致谢}
本研究得到了多方支持与帮助，在此表示衷心感谢。

感谢MIMIC-IV数据集的开发团队和Beth Israel Deaconess Medical Center，为研究社区提供了高质量的去标识化临床数据资源。感谢Microsoft GraphRAG团队开源了强大的知识图谱构建框架，为本研究的核心技术提供了基础。

特别感谢审稿人的宝贵意见和建设性建议，极大地提升了论文质量。

本研究由庆北国立大学提供计算资源支持。

\textbf{通讯作者：}YE SEVERIN（\texttt{6severin9@gmail.com}）

% =========================
% Appendix
% =========================
\section{Appendix}
本附录提供更详细的实验配置、提示模板、数据集构建说明与可复现性信息，以支持对本文结果的复核与复现。

\subsection{Appendix A: Experimental Setup}

\subsubsection{A.1 Hardware and Software}

所有实验在单机环境完成。本地推理主要使用单张 NVIDIA RTX 4090（24GB 显存），LoRA 微调使用单张 NVIDIA H200。服务器配备 Intel Xeon Platinum 8358（32 核）与 256GB 内存，存储为 2TB NVMe SSD。

软件环境如下：

\begin{itemize}
\tightlist
\item
  Python 3.10.12
\item
  PyTorch 2.1.0
\item
  Transformers 4.36.0
\item
  LangChain 1.0.0
\item
  GraphRAG 0.2.1
\item
  LanceDB 0.3.2
\end{itemize}

\subsubsection{A.2 Model Hyperparameters}

\textbf{Qwen 3 14b LoRA 微调：}

\begin{itemize}
\tightlist
\item
  学习率：2e-5
\item
  批大小：8（梯度累积步数 = 4）
\item
  训练轮数：3
\item
  预热步数：100
\item
  最大序列长度：2048
\item
  LoRA rank：16
\item
  LoRA alpha：32
\end{itemize}

\textbf{推理参数：}

\begin{itemize}
\tightlist
\item
  温度：0.3（检测），0.7（生成）
\item
  top-p：0.9
\item
  最大输出 tokens：512
\item
  停止序列：\texttt{{[}"\textbackslash{}n\textbackslash{}n\textbackslash{}n"{]}}
\end{itemize}

\subsection{Appendix B: Prompt Templates}

\subsubsection{B.1 Entity Extraction Prompt (Medical Customization)}

我们使用如下``协议式''指令模板，从患者的 EHR 数据中抽取实体与关系，用于构建患者特定知识图谱。

\textbf{Instruction（指令）}

你是一名医学知识图谱抽取系统。请从输入的 EHR 记录中抽取医学实体与关系，并以 JSON 输出。

\textbf{Constraints（约束）}

\begin{enumerate}
\def\labelenumi{\arabic{enumi}.}
\tightlist
\item
  \textbf{患者实体唯一性}

  \begin{itemize}
  \tightlist
  \item
    为整个记录创建且仅创建一个 \texttt{Patient} 实体。
  \item
    不要创建 \texttt{Patient\_1}、\texttt{Patient\_2} 或多个患者实体。
  \item
    所有诊断、用药、检查都必须链接到这一个 \texttt{Patient}。
  \end{itemize}
\item
  \textbf{化验检查归一化}

  \begin{itemize}
  \tightlist
  \item
    将单项指标归并到检查面板中，例如：

    \begin{itemize}
    \tightlist
    \item
      CBC：WBC、RBC、HGB、HCT、PLT
    \item
      肝功能：ALT、AST、ALK PHOS、BILIRUBIN
    \item
      肾功能：CREAT、UREA N、BUN
    \end{itemize}
  \item
    不要为每个指标创建单独的 \texttt{LAB\_TEST} 实体。
  \item
    为具体数值创建 \texttt{LAB\_RESULT} 实体，并通过关系指向对应 \texttt{LAB\_TEST}。
  \end{itemize}
\item
  \textbf{语言一致性}

  \begin{itemize}
  \tightlist
  \item
    实体名称统一使用中文。
  \item
    英文保留在描述字段中。
  \item
    示例：\texttt{name="高脂血症"}, \texttt{description="HYPERLIPIDEMIA"}。
  \end{itemize}
\end{enumerate}

\textbf{Entity Types（实体类型）}

\begin{itemize}
\tightlist
\item
  PATIENT, DIAGNOSIS, MEDICATION, LAB\_TEST, LAB\_RESULT
\item
  VITAL\_SIGN, SYMPTOM, PROCEDURE, DEPARTMENT
\end{itemize}

\textbf{Relation Types（关系类型）}

\begin{itemize}
\tightlist
\item
  HAS\_DIAGNOSIS, PRESCRIBED, SHOWS, UNDERWENT
\item
  HAS\_VITAL\_SIGN, TESTED\_BY, RESULT\_OF, TREATED\_IN
\end{itemize}

\textbf{Output Format（输出格式）}

输出 JSON，包含 \texttt{entities} 与 \texttt{relationships} 两个数组：

\begin{verbatim}
{
    "entities": [
        {
            "id": "...",
            "type": "PATIENT|DIAGNOSIS|...",
            "name": "...",
            "description": "..."
        }
    ],
    "relationships": [
        {
            "source": "...",
            "target": "...",
            "type": "HAS_DIAGNOSIS|PRESCRIBED|...",
            "description": "..."
        }
    ]
}
\end{verbatim}

\textbf{Why this design（设计动机）}

该模板通过``患者实体唯一性''和``化验面板归一化''约束，减少实体过度碎片化，并提高跨表证据对齐的一致性，从而为后续 GraphRAG 检索与事实核验提供更稳定的结构化支撑。

\subsubsection{B.2 Hallucination Detection Prompt}

我们使用如下模板对出院小结中的单句陈述进行核验：给定句子与从知识图谱检索到的 EHR 证据上下文，判断该句是否为幻觉，并输出结构化解释。

\textbf{Inputs（输入）}

\begin{itemize}
\tightlist
\item
  \texttt{sentence}: 待核验的句子
\item
  \texttt{evidence\_context}: 检索到的 EHR 证据上下文
\end{itemize}

\textbf{Task（任务）}

\begin{enumerate}
\def\labelenumi{\arabic{enumi}.}
\tightlist
\item
  判断句子是否被证据支持。
\item
  识别任何相互矛盾的信息。
\item
  分配证据等级：

  \begin{itemize}
  \tightlist
  \item
    E1：由结构化 EHR 字段支持
  \item
    E2：由非结构化文本支持
  \item
    E3：证据未提及，但根据上下文判断应视为有问题
  \item
    E4：与现有事实证据冲突（相互矛盾）
  \end{itemize}
\end{enumerate}

\textbf{Output Format（输出格式）}

输出 JSON：

\begin{verbatim}
{
    "sentence_index": 0,
    "sentence": "...",
    "hallucination_detected": false,
    "evidence_grade": "E1|E2|E3|E4",
    "explanation": "...",
    "supporting_evidence": ["..."],
    "conflicting_evidence": ["..."],
    "confidence_score": 0.0
}
\end{verbatim}

\textbf{Consistency Rule（一致性规则）}

若 \texttt{hallucination\_detected\ =\ true}，则 \texttt{evidence\_grade} 必须为 \texttt{E3} 或 \texttt{E4}。

\subsection{Appendix C: Error Analysis}

基于对典型误判案例的定性复核，主要错误模式如下。

\textbf{False Negatives（遗漏的 E4 幻觉）}

\begin{itemize}
\tightlist
\item
  医学同义词或缩写未能与 EHR 中的标准表述正确对齐。
\item
  需要跨时间点整合病程信息的陈述未被稳定识别为冲突。
\item
  生成句子依赖隐含临床含义，而非 EHR 中的显式记录。
\item
  图检索未返回足以支撑判别的关键证据。
\end{itemize}

\textbf{False Positives（正确陈述被标记）}

\begin{itemize}
\tightlist
\item
  生成文本对病历中的已有暗示作了保守外推，但原文并未显式逐字记录。
\item
  语义等价表达在字符串层面未完全对齐。
\item
  时间范围或事件先后关系的表述存在歧义。
\item
  图谱抽取阶段的实体或关系误差传导到最终判别。
\end{itemize}

\subsection{Appendix D: Dataset Details}

\subsubsection{D.1 Source Tables}

我们基于 Discharge-Me / MIMIC-IV 的多表字段构建实验数据。下述``表级记录规模''反映的是完整原始表视图的规模，用于说明数据来源；其规模不等同于本文按 train、valid、test\_phase\_1、test\_phase\_2 切分后的实验目录规模，也不应直接理解为本文 250 名患者实验子集的实际规模。

\textbf{Diagnosis Table（diagnosis.csv）}

\begin{itemize}
\tightlist
\item
  规模：143,285 条诊断记录（原始表视图）
\item
  关键字段：\texttt{subject\_id}, \texttt{stay\_id}, \texttt{icd\_code}, \texttt{icd\_version}, \texttt{icd\_title}
\item
  示例：\texttt{10000032}, \texttt{J18.9}, \texttt{ICD-10}, ``未特指病原体的肺炎''
\end{itemize}

\textbf{ED Stays and Triage Tables（edstays.csv / triage.csv）}

\begin{itemize}
\tightlist
\item
  用途：提供急诊就诊、分诊与生命体征信息
\item
  关键字段（标识与时间）：\texttt{subject\_id}, \texttt{stay\_id}, \texttt{intime}, \texttt{outtime}
\item
  关键字段（生命体征）：\texttt{temperature}, \texttt{heartrate}, \texttt{resprate}, \texttt{o2sat}, \texttt{sbp}, \texttt{dbp}, \texttt{pain}, \texttt{acuity}
\end{itemize}

\textbf{Discharge Target Table（discharge\_target.csv）}

\begin{itemize}
\tightlist
\item
  用途：提供 \texttt{brief\_hospital\_course} 与 \texttt{discharge\_instructions} 等目标 section 文本
\item
  关键字段：\texttt{subject\_id}, \texttt{note\_id}, \texttt{brief\_hospital\_course}, \texttt{discharge\_instructions}
\item
  说明：是生成与核验任务的核心文本来源之一
\end{itemize}

\textbf{Radiology Table（radiology.csv）}

\begin{itemize}
\tightlist
\item
  用途：补充放射学文本证据
\item
  关键字段：\texttt{subject\_id}, \texttt{stay\_id}, \texttt{text}, \texttt{charttime}
\item
  说明：为影像相关事实核验提供补充证据
\end{itemize}

\subsubsection{D.2 Data Preprocessing}

我们采用如下预处理流程将多源表字段组织为 GraphRAG 可用的输入格式：

\begin{enumerate}
\def\labelenumi{\arabic{enumi}.}
\tightlist
\item
  \textbf{多表连接}：在 \texttt{subject\_id} 与 \texttt{stay\_id} 上整合 diagnosis、discharge、discharge\_target、edstays、radiology、triage 六类来源。
\item
  \textbf{去重}：移除重复的 stay 记录（348 例）。
\item
  \textbf{缺失值处理}：

  \begin{itemize}
  \tightlist
  \item
    删除缺失出院文本的记录（1,203 例）。
  \item
    对 triage 中缺失的生命体征字段进行标准化处理。
  \end{itemize}
\item
  \textbf{文本清洗}：

  \begin{itemize}
  \tightlist
  \item
    移除 PHI 标记（例如 \texttt{{[}*\ *\ ...\ *\ *{]}}）。
  \item
    规范化空白字符。
  \item
    修正常见拼写错误。
  \end{itemize}
\item
  \textbf{格式转换}：CSV → JSONL，以兼容 GraphRAG。
\end{enumerate}

\subsection{Appendix E: Reproducibility Checklist}

\begin{itemize}
\tightlist
\item
  Code Availability：项目代码计划发布于 https://github.com/{[}username{]}/CuraView
\item
  Data Access：MIMIC-IV（需要 PhysioNet 凭证访问）
\item
  Model Checkpoints：HuggingFace Hub \texttt{{[}username{]}/curaview-qwen3-14b}
\item
  Training Configuration：详见附录 A.2
\item
  Random Seed：所有实验固定为 42
\item
  Hardware：LoRA 微调使用单张 NVIDIA H200；本地推理使用单张 RTX 4090（24GB）
\item
  Training Time：约 48 小时（LoRA 微调）
\item
  Dependencies：仓库 \texttt{requirements.txt}
\end{itemize}

% =========================
% References
% =========================
\bibliographystyle{unsrtnat}
\bibliography{ref}

\end{document}
